\section{Stelle ibride: EOS nucleonica e QCD}
\subsection{Regioni di densita' e vincoli teorici}

% 01
\begin{frame}{Capitolo 9: il problema delle stelle ibride}
\small
\begin{columns}[T,onlytextwidth]
\begin{column}{0.49\textwidth}
\textbf{Basse densita'}
\begin{itemize}
\item Descrizione con nucleoni, nuclei ed elettroni.
\item EOS nucleonica vincolata dall'interazione nucleare.
\item Vicino alla densita' di saturazione $n_0$.
\end{itemize}
\end{column}
\begin{column}{0.49\textwidth}
\textbf{Altissime densita'}
\begin{itemize}
\item Gradi di liberta': quark e gluoni.
\item QCD perturbativa, affidabile solo a potenziali chimici elevati.
\item EOS di materia di quark: $P_{\rm pQCD}(\mu_q)$.
\end{itemize}
\end{column}
\end{columns}
\medskip
La regione intermedia non e' fissata univocamente da questi due limiti.
\medskip
Scopo del capitolo: costruire EOS che li colleghino e studiarne le soluzioni TOV.
\end{frame}

% 02
\begin{frame}{Transizione adroni--quark: che cosa sappiamo?}
\small
\begin{columns}[T,onlytextwidth]
\begin{column}{0.49\textwidth}
Il passaggio fra gradi di liberta' adronici e quark puo' essere:
\begin{itemize}
\item un \emph{crossover} liscio;
\item una transizione continua con non analiticita' delle derivate;
\item una transizione del primo ordine con salto della densita' di energia.
\end{itemize}
\end{column}
\begin{column}{0.49\textwidth}
Le informazioni della QCD a $T$ elevata e $\mu_B$ piccolo non determinano la natura della transizione a
\[
T\simeq0,\qquad \mu_B\gg0.
\]
L'interpolazione fra le EOS non costituisce, da sola, prova di deconfinamento o di una transizione di fase.
\end{column}
\end{columns}
\end{frame}

% 03
\begin{frame}{Scala nucleare: densita' di saturazione}
\small
\begin{columns}[T,onlytextwidth]
\begin{column}{0.49\textwidth}
Per materia nucleare simmetrica a $T=0$:
\[
n_{\rm sat}\simeq 0.15\text{--}0.16\,\mathrm{fm}^{-3}.
\]
Massa media del nucleone:
\[
m_N\simeq 939\,\mathrm{MeV}.
\]
La scala di densita' di energia e'
\[
\varepsilon_0\simeq m_N n_{\rm sat}.
\]
\end{column}
\begin{column}{0.49\textwidth}
Usando $n_{\rm sat}=0.15\,\mathrm{fm}^{-3}$:
\[
\begin{aligned}
\varepsilon_0&\simeq 939\times0.15\\
&\simeq 141\,\mathrm{MeV\,fm}^{-3}.
\end{aligned}
\]
Il libro arrotonda a
\[
\varepsilon_0\simeq140\,\mathrm{MeV\,fm}^{-3}.
\]
Questa scala serve a esprimere le densita' dei capitoli 7--9.
\end{column}
\end{columns}
\end{frame}

% 04
\begin{frame}{Crosta e materia neutronica: gerarchia delle densita'}
\footnotesize
\begin{columns}[T,onlytextwidth]
\begin{column}{0.49\textwidth}
\textbf{Crosta esterna}
\[
\varepsilon\lesssim10^{-3}\varepsilon_0.
\]
Reticolo di nuclei ed elettroni degeneri; assenza di neutroni liberi.
\medskip
\textbf{Crosta interna}
\[
10^{-3}\varepsilon_0\lesssim\varepsilon
\lesssim\varepsilon_0/2.
\]
Reticolo di nuclei, elettroni e neutroni oltre il \emph{neutron drip}.
\end{column}
\begin{column}{0.49\textwidth}
\textbf{Materia nucleonica ricca di neutroni}
\[
\varepsilon\gtrsim\varepsilon_0/2.
\]
Si passa verso il fluido nucleonico del core.
\medskip
A densita' prossime e poco superiori a $\varepsilon_0$, il testo assume un'incertezza nucleonica che puo' gia' raggiungere circa il $20\%$.
\medskip
Le soglie indicate sono ordini di grandezza della descrizione utilizzata.
\end{column}
\end{columns}
\end{frame}

\subsection{Energia di simmetria: origine di $m_N-16+32$}

% 05
\begin{frame}{Composizione nucleonica e asimmetria isospinica}
\small
Definiamo
\[
n_B=n_n+n_p,\qquad Y_p=\frac{n_p}{n_B},
\]
\[
\delta\equiv\frac{n_n-n_p}{n_B}=1-2Y_p.
\]
\begin{columns}[T,onlytextwidth]
\begin{column}{0.49\textwidth}
\textbf{Materia simmetrica}
\[
n_n=n_p,\quad Y_p=\frac12,\quad\delta=0.
\]
Per simmetria di isospin, trascurando piccole rotture, $e(n_B,\delta)$ e' pari in $\delta$.
\end{column}
\begin{column}{0.49\textwidth}
\textbf{Materia neutronica pura}
\[
n_p=0,\quad Y_p=0,\quad\delta=1.
\]
La differenza energetica e' dominata dall'energia di simmetria, non dalla sola massa a riposo.
\end{column}
\end{columns}
\end{frame}

% 06
\begin{frame}{Derivazione dell'approssimazione parabolica}
\small
Sia $e(n_B,\delta)=\varepsilon(n_B,\delta)/n_B$ l'energia \emph{totale} per barione, inclusa la massa a riposo.
\medskip
Espandiamo a densita' fissata attorno a $\delta=0$:
\[
\begin{aligned}
e(n_B,\delta)={}&e(n_B,0)
+\left.\frac{\partial e}{\partial\delta}\right|_0\delta\\
&+\frac12\left.\frac{\partial^2e}{\partial\delta^2}\right|_0\delta^2
+O(\delta^4).
\end{aligned}
\]
La simmetria approssimata $n_n\leftrightarrow n_p$ implica $e(n_B,\delta)=e(n_B,-\delta)$: il termine lineare e' nullo.
\[
{e(n_B,\delta)\simeq e(n_B,0)+S(n_B)\delta^2},
\qquad
S(n_B)=\frac12\left.\frac{\partial^2e}{\partial\delta^2}\right|_0.
\]
\end{frame}

% 07
\begin{frame}{Da $m_N-16$ alla materia neutronica pura}
\small
\begin{columns}[T,onlytextwidth]
\begin{column}{0.49\textwidth}
\textbf{Input fenomenologici a saturazione}
\[
e(n_{\rm sat},0)\simeq m_N-16\,\mathrm{MeV},
\]
\[
S(n_{\rm sat})\equiv J\simeq32\,\mathrm{MeV}.
\]
$16$ MeV: legame per nucleone della materia simmetrica infinita.
\medskip
$32$ MeV: coefficiente di simmetria \emph{della materia nucleare} a saturazione.
\end{column}
\begin{column}{0.49\textwidth}
\textbf{Materia di neutroni puri}
\[
\delta=1\quad\Longrightarrow\quad \delta^2=1.
\]
\[
\begin{aligned}
e_{\rm PNM}(n_{\rm sat})
&\simeq e(n_{\rm sat},0)+J\\
&\simeq m_N-16+32\\
&=m_N+16\\
&\simeq955\,\mathrm{MeV}.
\end{aligned}
\]
\end{column}
\end{columns}
\medskip
\textit{Precisazione:} sono valori fenomenologici approssimati, non una previsione diretta della pQCD.
\end{frame}

% 08
\begin{frame}{Come nasce il termine di simmetria nelle masse nucleari?}
\small
Per un nucleo con $A$ nucleoni e $Z$ protoni, la formula semiempirica include
\[
E_b\supset a_vA-a_{\rm sym}^{(A)}\frac{(A-2Z)^2}{A}.
\]
Poiche' $\delta=(A-2Z)/A$, dividendo per $A$ e trascurando superficie, Coulomb e pairing:
\[
\frac{E_b}{A}\simeq a_v-a_{\rm sym}^{\rm bulk}\delta^2.
\]
L'energia totale per barione e'
\[
\begin{aligned}
e&\simeq m_N-\frac{E_b}{A}\\
&\simeq m_N-a_v+a_{\rm sym}^{\rm bulk}\delta^2.
\end{aligned}
\]
Nel limite della materia nucleare: $a_v\simeq16$ MeV e $a_{\rm sym}^{\rm bulk}=J\simeq32$ MeV.
\medskip
\textit{Nota:} il coefficiente di simmetria di un nucleo finito non coincide esattamente con $J$ perche' include effetti di superficie.
\end{frame}

% 09
\begin{frame}{Energia per barione e potenziale chimico non coincidono sempre}
\small
Per materia fredda con un solo numero barionico conservato:
\[
\mu_B=\frac{d\varepsilon}{dn_B},\qquad P=n_B\mu_B-\varepsilon,
\]
quindi
\[
{\mu_B=\frac{\varepsilon}{n_B}+\frac{P}{n_B}=e+\frac{P}{n_B}.}
\]
\begin{columns}[T,onlytextwidth]
\begin{column}{0.49\textwidth}
Materia \emph{simmetrica} alla saturazione:
\[
P(n_{\rm sat},\delta=0)=0
\]
e percio' $\mu_B=e\simeq m_N-16$ MeV.
\end{column}
\begin{column}{0.49\textwidth}
Materia \emph{neutronica} alla stessa densita':
\[
P_{\rm PNM}(n_{\rm sat})>0
\]
in generale, dunque $\mu_n>e_{\rm PNM}$.
\end{column}
\end{columns}
\medskip
Il valore $955$ MeV del testo fornisce la \emph{scala} $E/A$, non necessariamente il potenziale chimico esatto.
\end{frame}

% 10
\begin{frame}{Correzione esplicita: il ruolo della pendenza $L$}
\small
Scriviamo
\[
e(n_B,\delta)\simeq e_{\rm sat}(n_B)+S(n_B)\delta^2.
\]
Per neutroni puri a $n_{\rm sat}$, $\delta=1$, mentre $e_{\rm sat}'(n_{\rm sat})=0$.
\[
\begin{aligned}
P_{\rm PNM}(n_{\rm sat})
&=n_{\rm sat}^2\left.\frac{d e_{\rm PNM}}{dn_B}\right|_{n_{\rm sat}}\\
&\simeq n_{\rm sat}^2 S'(n_{\rm sat}).
\end{aligned}
\]
Definendo $L=3n_{\rm sat}S'(n_{\rm sat})$:
\[
P_{\rm PNM}(n_{\rm sat})\simeq\frac{n_{\rm sat}L}{3},
\qquad
{\mu_n\simeq m_N+16+\frac L3.}
\]
La stima del libro $\mu_B\sim955$ MeV trascura questo termine di pressione.
\end{frame}

\subsection{Vincolo perturbativo ad altissime densita'}

% 11
\begin{frame}{Limite pQCD: quanto e' distante dalla saturazione?}
\small
\begin{columns}[T,onlytextwidth]
\begin{column}{0.49\textwidth}
Il testo prende come scala indicativa:
\[
\mu_q\simeq900\,\mathrm{MeV},
\]
con
\[
\mu_B=3\mu_q\simeq2.7\,\mathrm{GeV}.
\]
All'estremo nucleonico:
\[
\mu_B\sim0.95\,\mathrm{GeV}.
\]
Differenza nei potenziali chimici: meno di un fattore $3$.
\end{column}
\begin{column}{0.49\textwidth}
Ma la densita' di energia cresce approssimativamente come
\[
\varepsilon_{\rm free}\propto\mu_q^4.
\]
Ne consegue un divario assai maggiore fra densita' nucleari e densita' a cui la pQCD e' controllata.
\medskip
La materia nel core delle stelle osservate ricade in gran parte nella regione intermedia.
\end{column}
\end{columns}
\end{frame}

% 12
\begin{frame}{Densita' del gas libero a $\mu_q=900$ MeV}
\small
Tre sapori massless, $T=0$:
\[
P_{\rm SB}=\frac{3\mu_q^4}{4\pi^2},
\qquad \varepsilon_{\rm SB}=3P_{\rm SB}=\frac{9\mu_q^4}{4\pi^2}.
\]
Ripristinando $\hbar c$:
\[
\varepsilon_{\rm SB}[\mathrm{MeV\,fm^{-3}}]
=\frac{9(900\,\mathrm{MeV})^4}{4\pi^2(\hbar c)^3},
\quad \hbar c\simeq197.327\,\mathrm{MeV\,fm}.
\]
Numericalmente:
\[
\varepsilon_{\rm SB}\simeq1.95\times10^4\,\mathrm{MeV\,fm^{-3}}
\simeq140\,\varepsilon_0.
\]
Con le correzioni interagenti, il libro stima $\varepsilon_{\rm pQCD}\sim70\text{--}100\,\varepsilon_0$.
\end{frame}

\subsection{Figura 9.1: il divario fra le EOS}

% 13
\begin{frame}{Figura 9.1: EOS nucleonica, regione incerta e pQCD}
\small
\begin{columns}[T,onlytextwidth]
\begin{column}{0.63\textwidth}
    \begin{figure}[!ht]\includegraphics[trim={0cm 0cm 0 0},clip,keepaspectratio,width=0.95\textwidth]{Schaffner/EOS-crust-pQCD}\label{fig:Schaffner_EOS-crust-pQCD}\end{figure}
\end{column}
\begin{column}{0.35\textwidth}
\textbf{Asse verticale}
\[
P\ [\mathrm{MeV\,fm^{-3}}].
\]
\textbf{Asse orizzontale}
\[
\widetilde\mu_q=\mu_q-\frac{\mu_{\rm Fe}}3.
\]
Entrambi gli assi: scala logaritmica.
\medskip
A sinistra: croste e materia nucleonica. A destra: pQCD.
\end{column}
\end{columns}
\end{frame}

% 14
\begin{frame}{Perche' traslare il potenziale chimico dei quark?}
\small
Il potenziale chimico di riferimento del ferro e' assunto come
\[
\mu_{\rm Fe}\simeq930\,\mathrm{MeV},\qquad \frac{\mu_{\rm Fe}}3\simeq310\,\mathrm{MeV}.
\]
Definendo $\widetilde\mu_q=\mu_q-\mu_{\rm Fe}/3$,
\[
\mu_q=310\,\mathrm{MeV}+\widetilde\mu_q,
\qquad \mu_B=930\,\mathrm{MeV}+3\widetilde\mu_q.
\]
Esempio:
\[
\widetilde\mu_q=10\,\mathrm{MeV}
\quad\Longrightarrow\quad\mu_B\simeq960\,\mathrm{MeV}.
\]
La traslazione rende leggibili su scala logaritmica piccoli incrementi di potenziale chimico rispetto alla materia nucleare stabile.
\end{frame}

% 15
\begin{frame}{Lettura dei confini nella Figura 9.1}
\small
Le linee verticali delimitano indicativamente:
\[
\begin{array}{ll}
\widetilde\mu_q\simeq3\,\mathrm{MeV} & \text{fine crosta esterna},\\[2mm]
\widetilde\mu_q\simeq7.4\,\mathrm{MeV} & \text{fine crosta interna},\\[2mm]
\widetilde\mu_q\simeq16\,\mathrm{MeV} & \text{estremo nucleonico}.
\end{array}
\]
Il margine inferiore pQCD si trova molto piu' a destra ($\widetilde\mu_q\sim560$ MeV nel grafico).
\medskip
La zona centrale, con il punto interrogativo, richiede una EOS interpolata; non equivale a una regione sperimentalmente accertata di quark liberi.
\medskip
La fascia vicina all'asse orizzontale segnala che le densita' centrali delle stelle di massa massima ricadono proprio nell'intervallo incerto.
\end{frame}

\section{Interpolazione con politropiche a tratti}
\subsection{Termodinamica a temperatura zero}

% 16
\begin{frame}{Politropiche a tratti: definizione}
\small
Per ogni intervallo di densita' si assume
\[
{P(n)=Kn^{\Gamma}},
\]
con $n=n_B$, $K$ costante e $\Gamma$ indice politropico.
\[
\Gamma=\frac{d\ln P}{d\ln n}.
\]
In una EOS a tratti:
\[
P(n)=\begin{cases}
K_1n^{\Gamma_1},&n_0\leq n<n_1,\\
K_2n^{\Gamma_2},&n_1\leq n<n_2,\\
\cdots&\cdots
\end{cases}
\]
Le costanti di ciascun segmento devono essere determinate con condizioni di raccordo termodinamico.
\end{frame}

% 17
\begin{frame}{Identita' termodinamiche essenziali}
\small
A temperatura zero, materia neutra in equilibrio con un'unica carica conservata indipendente:
\[
d\varepsilon=\mu\,dn,\qquad
P=\mu n-\varepsilon,\qquad dP=n\,d\mu.
\]
Da esse seguono
\[
{\mu=\frac{d\varepsilon}{dn}},
\qquad
{n=\frac{dP}{d\mu}}.
\]
\begin{columns}[T,onlytextwidth]
\begin{column}{0.49\textwidth}
La forma $P(n)$ \emph{non basta} a definire autonomamente l'energia: occorre integrare l'identita' termodinamica.
\end{column}
\begin{column}{0.49\textwidth}
La continuita' di $P(\mu)$ e della sua derivata garantisce un raccordo senza salto della densita' barionica.
\end{column}
\end{columns}
\end{frame}

% 18
\begin{frame}{Derivazione dell'identita' (9.2)}
\small
Partiamo da
\[
P=n\mu-\varepsilon,
\qquad\mu=\frac{d\varepsilon}{dn}.
\]
Sostituiamo il potenziale chimico:
\[
P=n\frac{d\varepsilon}{dn}-\varepsilon.
\]
Dividendo per $n^2$:
\[
\frac P{n^2}=\frac1n\frac{d\varepsilon}{dn}-\frac{\varepsilon}{n^2}.
\]
Il membro destro e' la derivata di un rapporto:
\[
{P=n^2\frac{d}{dn}\left(\frac{\varepsilon}{n}\right).}
\]
\end{frame}

% 19
\begin{frame}{Dalla politropica alla densita' di energia (9.3)}
\small
Assumiamo in un intervallo
\[
P(n)=Kn^\Gamma.
\]
Dall'identita' precedente:
\[
Kn^\Gamma=n^2\frac{d}{dn}\left(\frac{\varepsilon}{n}\right).
\]
Dividiamo per $n^2$:
\[
{\frac{d}{dn}\left(\frac{\varepsilon}{n}\right)=Kn^{\Gamma-2}.}
\]
Per $\Gamma\ne1$, integriamo:
\[
\frac{\varepsilon}{n}=K\int n^{\Gamma-2}dn+c
=\frac{K}{\Gamma-1}n^{\Gamma-1}+c.
\]
\end{frame}

% 20
\begin{frame}{Densita' di energia politropica (9.4)}
\small
Moltiplicando il risultato integrato per $n$:
\[
{\varepsilon(n)=\frac{K}{\Gamma-1}n^\Gamma+cn.}
\]
Poiche' $P=Kn^\Gamma$,
\[
{\varepsilon(n)=\frac{P(n)}{\Gamma-1}+cn.}
\]
\begin{columns}[T,onlytextwidth]
\begin{column}{0.49\textwidth}
Il primo termine e' determinato dalla legge politropica $P(n)$.
\end{column}
\begin{column}{0.49\textwidth}
Il termine $cn$ integra l'informazione sull'energia per barione al raccordo; in generale $c\ne m_N$.
\end{column}
\end{columns}
\end{frame}

% 21
\begin{frame}{Potenziale chimico della politropica (9.5)}
\small
Per definizione:
\[
\mu(n)=\frac{d\varepsilon(n)}{dn}.
\]
Derivando termine a termine:
\[
\begin{aligned}
\mu(n)&=\frac{d}{dn}\left[\frac{K}{\Gamma-1}n^\Gamma+cn\right]\\
&=\frac{K\Gamma}{\Gamma-1}n^{\Gamma-1}+c.
\end{aligned}
\]
Quindi
\[
{\mu(n)=c+\frac{K\Gamma}{\Gamma-1}n^{\Gamma-1}.}
\]
Verifica: $n\mu-\varepsilon=Kn^\Gamma=P$.
\end{frame}

\subsection{Raccordo al punto $(n_0,P_0,\mu_0)$}

% 22
\begin{frame}{Primo vincolo di raccordo: pressione}
\small
Sia $(n_0,P_0,\mu_0)$ il punto terminale dell'EOS nota.
\medskip
Per il segmento politropico imponiamo
\[
P(n_0)=P_0.
\]
Con $P(n)=Kn^\Gamma$:
\[
P_0=Kn_0^\Gamma
\quad\Longrightarrow\quad
{K=\frac{P_0}{n_0^\Gamma}.}
\]
\begin{columns}[T,onlytextwidth]
\begin{column}{0.49\textwidth}
$\Gamma$ determina la forma del segmento.
\end{column}
\begin{column}{0.49\textwidth}
$K$ non e' libero una volta scelta la pressione al raccordo.
\end{column}
\end{columns}
\medskip
Il simbolo $n_0$ qui identifica il \emph{punto di raccordo}, non necessariamente la densita' di saturazione nucleare.
\end{frame}

% 23
\begin{frame}{Secondo vincolo: continuita' del potenziale (9.6)}
\small
Partiamo da
\[
\mu(n)=c+\frac{K\Gamma}{\Gamma-1}n^{\Gamma-1}.
\]
Imponiamo $\mu(n_0)=\mu_0$:
\[
\mu_0=c+\frac{K\Gamma}{\Gamma-1}n_0^{\Gamma-1}.
\]
Si ricava
\[
{c=\mu_0-\frac{K\Gamma}{\Gamma-1}n_0^{\Gamma-1}.}
\]
Inserendo $K=P_0/n_0^\Gamma$:
\[
{c=\mu_0-\frac{\Gamma}{\Gamma-1}\frac{P_0}{n_0}.}
\]
Le due condizioni fissano $K$ e $c$; rimane da scegliere $\Gamma$.
\end{frame}

% 24
\begin{frame}{Controllo del raccordo: anche $\varepsilon$ e' continua}
\small
Alla densita' $n_0$,
\[
\varepsilon(n_0)=\frac{P_0}{\Gamma-1}+cn_0.
\]
Usiamo
\[
cn_0=\mu_0n_0-\frac{\Gamma}{\Gamma-1}P_0.
\]
Segue
\[
\begin{aligned}
\varepsilon(n_0)
&=\frac{P_0}{\Gamma-1}+\mu_0n_0
-\frac{\Gamma P_0}{\Gamma-1}\\
&=\mu_0n_0-P_0.
\end{aligned}
\]
\[
{\varepsilon_0+P_0=\mu_0n_0.}
\]
La densita' di energia coincide con quella dell'EOS nota se il punto soddisfa la stessa identita' termodinamica.
\end{frame}

\subsection{Da $P(n)$ a $P(\mu)$ e $\varepsilon(P)$}

% 25
\begin{frame}{Invertire il potenziale chimico: primo passaggio}
\small
Partiamo dalla (9.5):
\[
\mu-c=\frac{K\Gamma}{\Gamma-1}n^{\Gamma-1}.
\]
Isolando la potenza di $n$:
\[
n^{\Gamma-1}=\frac{\Gamma-1}{K\Gamma}(\mu-c).
\]
Elevando a $1/(\Gamma-1)$:
\[
{n(\mu)=\left[\frac{\Gamma-1}{K\Gamma}(\mu-c)\right]^{1/(\Gamma-1)}.}
\]
L'argomento della potenza deve rimanere positivo nel tratto fisico scelto.
\end{frame}

% 26
\begin{frame}{Densita' in funzione di $\mu$: raccordo (9.7)}
\small
Sostituiamo
\[
c=\mu_0-\frac{K\Gamma}{\Gamma-1}n_0^{\Gamma-1}
\]
nel termine $\mu-c$:
\[
\mu-c=(\mu-\mu_0)+\frac{K\Gamma}{\Gamma-1}n_0^{\Gamma-1}.
\]
Pertanto
\[
{
 n(\mu)=\left[n_0^{\Gamma-1}
 +\frac{\Gamma-1}{K\Gamma}(\mu-\mu_0)
 \right]^{1/(\Gamma-1)}.}
\]
Controllo immediato: $n(\mu_0)=n_0$.
\end{frame}

% 27
\begin{frame}{Pressione in funzione di $\mu$: Eq. (9.8)}
\small
Partendo da $P=Kn^\Gamma$ e dalla (9.7):
\[
{
P(\mu)=K\left[
 n_0^{\Gamma-1}+\frac{\Gamma-1}{K\Gamma}(\mu-\mu_0)
\right]^{\Gamma/(\Gamma-1)}.}
\]
Al raccordo $\mu=\mu_0$:
\[
P(\mu_0)=Kn_0^\Gamma=P_0.
\]
La funzione pressione e' quindi continua. Rimane da verificarne la derivata.
\end{frame}

% 28
\begin{frame}{Verifica: $dP/d\mu=n$ in ogni segmento}
\small
Definiamo
\[
X(\mu)=n_0^{\Gamma-1}+\frac{\Gamma-1}{K\Gamma}(\mu-\mu_0).
\]
Allora $P(\mu)=KX^{\Gamma/(\Gamma-1)}$ e
\[
\begin{aligned}
\frac{dP}{d\mu}
&=K\frac{\Gamma}{\Gamma-1}X^{1/(\Gamma-1)}
\frac{\Gamma-1}{K\Gamma}\\
&=X^{1/(\Gamma-1)}=n(\mu).
\end{aligned}
\]
Per $\mu=\mu_0$:
\[
{\left.\frac{dP}{d\mu}\right|_{\mu_0}=n_0.}
\]
Questa e' la condizione sulla pendenza usata per il raccordo liscio delle politropiche.
\end{frame}

% 29
\begin{frame}{Dalla pressione alla densita' per le TOV}
\small
Per integrare le TOV occorre $\varepsilon(P)$.
\medskip
Da $P=Kn^\Gamma$:
\[
n=\left(\frac PK\right)^{1/\Gamma}.
\]
Poiche' $P_0=Kn_0^\Gamma$, equivalentemente
\[
{n(P)=n_0\left(\frac P{P_0}\right)^{1/\Gamma}.}
\]
Usiamo la densita' di energia gia' integrata:
\[
\varepsilon(n)=\frac{P(n)}{\Gamma-1}+cn.
\]
Sostituiremo $n(P)$ e poi il valore di $c$ fissato al raccordo.
\end{frame}

% 30
\begin{frame}{EOS per le TOV: derivazione di (9.9)}
\small
Sostituendo $n(P)$:
\[
\varepsilon(P)=\frac{P}{\Gamma-1}+cn_0\left(\frac{P}{P_0}\right)^{1/\Gamma}.
\]
Dal vincolo sul potenziale:
\[
cn_0=\mu_0n_0-\frac{\Gamma}{\Gamma-1}P_0.
\]
Si ottiene
\[
{
\begin{aligned}
\varepsilon(P)={}&\frac{P}{\Gamma-1}\\
&+\left(\mu_0n_0-\frac{\Gamma}{\Gamma-1}P_0\right)
\left(\frac{P}{P_0}\right)^{1/\Gamma}.
\end{aligned}}
\]
Il termine in $P^{1/\Gamma}$ conserva l'informazione termodinamica del raccordo.
\end{frame}

% 31
\begin{frame}{Il caso speciale $\Gamma=1$}
\small
Le formule precedenti con $1/(\Gamma-1)$ richiedono $\Gamma\ne1$.
\medskip
Se $P=Kn$, l'equazione (9.2) diventa
\[
\frac{d}{dn}\left(\frac{\varepsilon}{n}\right)=\frac Kn.
\]
Integrando, con densita' di riferimento $n_{\rm ref}$,
\[
{\varepsilon(n)=Kn\ln\frac{n}{n_{\rm ref}}+cn.}
\]
Derivando:
\[
{\mu(n)=K\left[\ln\frac{n}{n_{\rm ref}}+1\right]+c.}
\]
Il segmento con $\Gamma=1$ si tratta dunque separatamente mediante un logaritmo.
\end{frame}

% 32
\begin{frame}{Algoritmo di raccordo per piu' politropiche}
\small
\begin{enumerate}
\item Partire dal punto nucleonico $(n_0,P_0,\mu_0)$.
\item Scegliere $\Gamma_1$ e determinare $K_1,c_1$ imponendo $P_0,\mu_0$.
\item Proseguire fino a un nuovo punto $(n_1,P_1,\mu_1)$.
\item Scegliere $\Gamma_2$ e imporre le stesse condizioni per il secondo segmento.
\item Raccordare all'EOS pQCD, controllando vincoli fisici lungo tutta la curva.
\end{enumerate}
\medskip
Al raccordo liscio:
\[
P_-(\mu_i)=P_+(\mu_i),\qquad
\left.\frac{dP_-}{d\mu}\right|_{\mu_i}
=\left.\frac{dP_+}{d\mu}\right|_{\mu_i}=n_i.
\]
\end{frame}

\subsection{Causalita', rigidita' e transizioni di fase}

% 33
\begin{frame}{Derivazione della velocita' del suono}
\small
A $T=0$, in un fluido barotropico:
\[
c_s^2=\frac{dP}{d\varepsilon}
=\frac{dP/dn}{d\varepsilon/dn}
\qquad(c=1).
\]
Per $P=Kn^\Gamma$:
\[
\frac{dP}{dn}=\Gamma Kn^{\Gamma-1}=\frac{\Gamma P}{n}.
\]
Dalla termodinamica:
\[
\frac{d\varepsilon}{dn}=\mu=\frac{\varepsilon+P}{n}.
\]
Quindi
\[
{c_s^2=\frac{\Gamma P}{\varepsilon+P}.}
\]
Causalita' e stabilita' locale richiedono $0\leq c_s^2\leq1$.
\end{frame}

% 34
\begin{frame}{Raccordo liscio e transizione del primo ordine}
\small
\begin{columns}[T,onlytextwidth]
\begin{column}{0.49\textwidth}
\textbf{Interpolazione liscia}
\[
P_H(\mu_i)=P_Q(\mu_i),
\]
\[
n_H(\mu_i)=n_Q(\mu_i).
\]
Per $\varepsilon=\mu n-P$ risulta anche $\varepsilon_H=\varepsilon_Q$ al punto.
\medskip
Nessun salto di densita' introdotto dal raccordo.
\end{column}
\begin{column}{0.49\textwidth}
\textbf{Primo ordine (Maxwell)}
\[
P_H(\mu_t)=P_Q(\mu_t),
\]
\[
n_Q\ne n_H.
\]
A pressione e potenziale comuni,
\[
{\Delta\varepsilon=\mu_t\Delta n_B.}
\]
La pressione continua, ma la derivata $dP/d\mu$ salta.
\end{column}
\end{columns}
\end{frame}

% 35
\begin{frame}{Quali politropiche sono fisicamente ammissibili?}
\small
\begin{columns}[T,onlytextwidth]
\begin{column}{0.49\textwidth}
\textbf{Vincoli locali}
\begin{itemize}
\item $n=dP/d\mu>0$;
\item stabilita' meccanica $dP/dn>0$;
\item causalita' $c_s^2=dP/d\varepsilon\leq1$.
\end{itemize}
\end{column}
\begin{column}{0.49\textwidth}
\textbf{Vincolo globale}
\[
M_{\max}\gtrsim2M_\odot
\]
per essere compatibili con l'esistenza di stelle massicce osservate.
\medskip
Una EOS troppo soffice non sostiene grandi masse; una troppo rigida puo' violare la causalita'.
\end{column}
\end{columns}
\medskip
La rigidita' stellare si giudica da $P(\varepsilon)$ e dalla velocita' del suono, non dal solo valore di $\Gamma$ isolatamente.
\end{frame}

\section{Massa--raggio delle stelle ibride}
\subsection{Figura 9.2 e vincoli astrofisici}

% 36
\begin{frame}{Figura 9.2: banda delle relazioni massa--raggio}
\small
\begin{columns}[T,onlytextwidth]
\begin{column}{0.64\textwidth}
    \begin{figure}[!ht]\includegraphics[trim={0cm 0cm 0 0},clip,keepaspectratio,width=0.95\textwidth]{Schaffner/MR-NS-hybrid}\label{fig:Schaffner_MR-NS-hybrid}\end{figure}
\end{column}
\begin{column}{0.34\textwidth}
Banda grigia: famiglie di EOS interpolate fra limite nucleonico e pQCD.
\medskip
Curve I, II e III: tre politropiche rappresentative.
\medskip
Linee: buchi neri, Buchdahl, limite causale e compattezze costanti.
\end{column}
\end{columns}
\end{frame}

% 37
\begin{frame}{Tre sequenze rappresentative nella Figura 9.2}
\footnotesize
\begin{columns}[T,onlytextwidth]
\begin{column}{0.33\textwidth}
\textbf{I: piu' compatta}
\medskip
Raggi relativamente piccoli.
\[
\varepsilon_{c,\max}\sim10\varepsilon_0.
\]
La configurazione di massa massima raggiunge densita' centrali elevate.
\end{column}
\begin{column}{0.33\textwidth}
\textbf{II: massa massima maggiore}
\[
M_{\max}\simeq2.4M_\odot,
\]
\[
\varepsilon_{c,\max}\sim5\varepsilon_0.
\]
Mostra come una EOS adeguatamente rigida possa sostenere masse elevate.
\end{column}
\begin{column}{0.33\textwidth}
\textbf{III: raggi maggiori}
\medskip
Sequenza piu' a destra nel diagramma.
\[
\varepsilon_{c,\max}\sim5\varepsilon_0.
\]
Raggi piu' ampi alle masse considerate.
\end{column}
\end{columns}
\end{frame}

% 38
\begin{frame}{Densita' centrali: cosa sondano le osservazioni?}
\small
Per la famiglia di EOS del testo:
\[
\begin{array}{c|c|c}
M & R & \varepsilon_c/\varepsilon_0\\ \hline
1.4M_\odot & 11\text{--}14.5\,\mathrm{km} & 1.6\text{--}4\\[1mm]
2.0M_\odot & 10\text{--}15\,\mathrm{km} & 2\text{--}8
\end{array}
\]
Confronto fra scale:
\[
\varepsilon_c\sim\text{pochi }\varepsilon_0,
\qquad \varepsilon_{\rm pQCD}\gtrsim70\,\varepsilon_0.
\]
Le osservazioni delle stelle compatte mettono alla prova la regione intermedia dell'EOS, non direttamente il regime di pQCD asintotica.
\medskip
{\scriptsize Intervalli del lavoro citato nel libro; non sono vincoli osservativi aggiornati.}
\end{frame}

% 39
\begin{frame}{Compattezza, causalita' e massa osservata}
\small
\begin{columns}[T,onlytextwidth]
\begin{column}{0.49\textwidth}
\textbf{Compattezza}
\[
\mathcal C=\frac{GM}{Rc^2}.
\]
Rette di riferimento nel grafico:
\[
\mathcal C=0.1,\ 0.2,\ 0.3.
\]
Limiti esaminati nel capitolo 8:
\[
\mathcal C_{\rm BH}=\tfrac12,\quad
\mathcal C_{\rm Buchdahl}=\tfrac49.
\]
\end{column}
\begin{column}{0.49\textwidth}
\textbf{Vincoli sull'EOS}
\[
c_s^2\leq1,\qquad M_{\max}\geq2M_\odot.
\]
La causalita' impedisce una EOS arbitrariamente rigida.
\medskip
La massa massima osservata impedisce una EOS arbitrariamente soffice.
\medskip
La loro combinazione restringe la banda ammissibile nel diagramma $M$--$R$.
\end{column}
\end{columns}
\end{frame}

\subsection{Riepilogo}

% 40
\begin{frame}{Risultato del \S9.1: interpolare senza inventare la microfisica}
\small
\begin{enumerate}
\item EOS nucleonica vincolata vicino a $\varepsilon_0$.
\item EOS pQCD controllata solo intorno a $\varepsilon\gtrsim70\text{--}100\varepsilon_0$.
\item Politropiche a tratti raccordate con continuita' di $P$, $\mu$ e $n$.
\item Verifica di causalita', stabilita' e massa massima.
\item Integrazione delle TOV per ottenere una banda di curve $M(R)$.
\end{enumerate}
\medskip
Il raccordo liscio e' una \emph{parametrizzazione} dell'EOS intermedia, non una dimostrazione che la transizione adroni--quark sia un crossover.
\end{frame}

% 41
\begin{frame}{Formule da ricordare: politropiche termodinamiche}
\footnotesize
\begin{columns}[T,onlytextwidth]
\begin{column}{0.49\textwidth}
\[
P(n)=Kn^\Gamma,
\]
\[
\varepsilon(n)=\frac{Kn^\Gamma}{\Gamma-1}+cn,
\]
\[
\mu(n)=\frac{K\Gamma}{\Gamma-1}n^{\Gamma-1}+c,
\]
\[
c=\mu_0-\frac{\Gamma P_0}{(\Gamma-1)n_0}.
\]
\end{column}
\begin{column}{0.49\textwidth}
\[
n=\frac{dP}{d\mu},
\]
\[
\varepsilon(P)=\frac{P}{\Gamma-1}
+\left(\mu_0n_0-\frac{\Gamma P_0}{\Gamma-1}\right)
\left(\frac{P}{P_0}\right)^{1/\Gamma},
\]
\[
c_s^2=\frac{\Gamma P}{\varepsilon+P}\leq1.
\]
\end{column}
\end{columns}
\medskip
\textbf{Energia di simmetria:} $e_{\rm PNM}(n_{\rm sat})\simeq m_N-16+32=m_N+16\simeq955$ MeV.
\end{frame}

\section{Stelle ibride: EOS nucleonica e QCD}
\subsection{Regioni di densita' e vincoli teorici}

% 01
\begin{frame}{Capitolo 9: il problema delle stelle ibride}
\small
\begin{columns}[T,onlytextwidth]
\begin{column}{0.49\textwidth}
\textbf{Basse densita'}
\begin{itemize}
\item Descrizione con nucleoni, nuclei ed elettroni.
\item EOS nucleonica vincolata dall'interazione nucleare.
\item Vicino alla densita' di saturazione $n_0$.
\end{itemize}
\end{column}
\begin{column}{0.49\textwidth}
\textbf{Altissime densita'}
\begin{itemize}
\item Gradi di liberta': quark e gluoni.
\item QCD perturbativa, affidabile solo a potenziali chimici elevati.
\item EOS di materia di quark: $P_{\rm pQCD}(\mu_q)$.
\end{itemize}
\end{column}
\end{columns}
\medskip
La regione intermedia non e' fissata univocamente da questi due limiti.
\medskip
Scopo del capitolo: costruire EOS che li colleghino e studiarne le soluzioni TOV.
\end{frame}

% 02
\begin{frame}{Transizione adroni--quark: che cosa sappiamo?}
\small
\begin{columns}[T,onlytextwidth]
\begin{column}{0.49\textwidth}
Il passaggio fra gradi di liberta' adronici e quark puo' essere:
\begin{itemize}
\item un \emph{crossover} liscio;
\item una transizione continua con non analiticita' delle derivate;
\item una transizione del primo ordine con salto della densita' di energia.
\end{itemize}
\end{column}
\begin{column}{0.49\textwidth}
Le informazioni della QCD a $T$ elevata e $\mu_B$ piccolo non determinano la natura della transizione a
\[
T\simeq0,\qquad \mu_B\gg0.
\]
L'interpolazione fra le EOS non costituisce, da sola, prova di deconfinamento o di una transizione di fase.
\end{column}
\end{columns}
\end{frame}

% 03
\begin{frame}{Scala nucleare: densita' di saturazione}
\small
\begin{columns}[T,onlytextwidth]
\begin{column}{0.49\textwidth}
Per materia nucleare simmetrica a $T=0$:
\[
n_{\rm sat}\simeq 0.15\text{--}0.16\,\mathrm{fm}^{-3}.
\]
Massa media del nucleone:
\[
m_N\simeq 939\,\mathrm{MeV}.
\]
La scala di densita' di energia e'
\[
\varepsilon_0\simeq m_N n_{\rm sat}.
\]
\end{column}
\begin{column}{0.49\textwidth}
Usando $n_{\rm sat}=0.15\,\mathrm{fm}^{-3}$:
\[
\begin{aligned}
\varepsilon_0&\simeq 939\times0.15\\
&\simeq 141\,\mathrm{MeV\,fm}^{-3}.
\end{aligned}
\]
Il libro arrotonda a
\[
\varepsilon_0\simeq140\,\mathrm{MeV\,fm}^{-3}.
\]
Questa scala serve a esprimere le densita' dei capitoli 7--9.
\end{column}
\end{columns}
\end{frame}

% 04
\begin{frame}{Crosta e materia neutronica: gerarchia delle densita'}
\footnotesize
\begin{columns}[T,onlytextwidth]
\begin{column}{0.49\textwidth}
\textbf{Crosta esterna}
\[
\varepsilon\lesssim10^{-3}\varepsilon_0.
\]
Reticolo di nuclei ed elettroni degeneri; assenza di neutroni liberi.
\medskip
\textbf{Crosta interna}
\[
10^{-3}\varepsilon_0\lesssim\varepsilon
\lesssim\varepsilon_0/2.
\]
Reticolo di nuclei, elettroni e neutroni oltre il \emph{neutron drip}.
\end{column}
\begin{column}{0.49\textwidth}
\textbf{Materia nucleonica ricca di neutroni}
\[
\varepsilon\gtrsim\varepsilon_0/2.
\]
Si passa verso il fluido nucleonico del core.
\medskip
A densita' prossime e poco superiori a $\varepsilon_0$, il testo assume un'incertezza nucleonica che puo' gia' raggiungere circa il $20\%$.
\medskip
Le soglie indicate sono ordini di grandezza della descrizione utilizzata.
\end{column}
\end{columns}
\end{frame}

\subsection{Energia di simmetria: origine di $m_N-16+32$}

% 05
\begin{frame}{Composizione nucleonica e asimmetria isospinica}
\small
Definiamo
\[
n_B=n_n+n_p,\qquad Y_p=\frac{n_p}{n_B},
\]
\[
\delta\equiv\frac{n_n-n_p}{n_B}=1-2Y_p.
\]
\begin{columns}[T,onlytextwidth]
\begin{column}{0.49\textwidth}
\textbf{Materia simmetrica}
\[
n_n=n_p,\quad Y_p=\frac12,\quad\delta=0.
\]
Per simmetria di isospin, trascurando piccole rotture, $e(n_B,\delta)$ e' pari in $\delta$.
\end{column}
\begin{column}{0.49\textwidth}
\textbf{Materia neutronica pura}
\[
n_p=0,\quad Y_p=0,\quad\delta=1.
\]
La differenza energetica e' dominata dall'energia di simmetria, non dalla sola massa a riposo.
\end{column}
\end{columns}
\end{frame}

% 06
\begin{frame}{Derivazione dell'approssimazione parabolica}
\small
Sia $e(n_B,\delta)=\varepsilon(n_B,\delta)/n_B$ l'energia \emph{totale} per barione, inclusa la massa a riposo.
\medskip
Espandiamo a densita' fissata attorno a $\delta=0$:
\[
\begin{aligned}
e(n_B,\delta)={}&e(n_B,0)
+\left.\frac{\partial e}{\partial\delta}\right|_0\delta\\
&+\frac12\left.\frac{\partial^2e}{\partial\delta^2}\right|_0\delta^2
+O(\delta^4).
\end{aligned}
\]
La simmetria approssimata $n_n\leftrightarrow n_p$ implica $e(n_B,\delta)=e(n_B,-\delta)$: il termine lineare e' nullo.
\[
{e(n_B,\delta)\simeq e(n_B,0)+S(n_B)\delta^2},
\qquad
S(n_B)=\frac12\left.\frac{\partial^2e}{\partial\delta^2}\right|_0.
\]
\end{frame}

% 07
\begin{frame}{Da $m_N-16$ alla materia neutronica pura}
\small
\begin{columns}[T,onlytextwidth]
\begin{column}{0.49\textwidth}
\textbf{Input fenomenologici a saturazione}
\[
e(n_{\rm sat},0)\simeq m_N-16\,\mathrm{MeV},
\]
\[
S(n_{\rm sat})\equiv J\simeq32\,\mathrm{MeV}.
\]
$16$ MeV: legame per nucleone della materia simmetrica infinita.
\medskip
$32$ MeV: coefficiente di simmetria \emph{della materia nucleare} a saturazione.
\end{column}
\begin{column}{0.49\textwidth}
\textbf{Materia di neutroni puri}
\[
\delta=1\quad\Longrightarrow\quad \delta^2=1.
\]
\[
\begin{aligned}
e_{\rm PNM}(n_{\rm sat})
&\simeq e(n_{\rm sat},0)+J\\
&\simeq m_N-16+32\\
&=m_N+16\\
&\simeq955\,\mathrm{MeV}.
\end{aligned}
\]
\end{column}
\end{columns}
\medskip
\textit{Precisazione:} sono valori fenomenologici approssimati, non una previsione diretta della pQCD.
\end{frame}

% 08
\begin{frame}{Come nasce il termine di simmetria nelle masse nucleari?}
\small
Per un nucleo con $A$ nucleoni e $Z$ protoni, la formula semiempirica include
\[
E_b\supset a_vA-a_{\rm sym}^{(A)}\frac{(A-2Z)^2}{A}.
\]
Poiche' $\delta=(A-2Z)/A$, dividendo per $A$ e trascurando superficie, Coulomb e pairing:
\[
\frac{E_b}{A}\simeq a_v-a_{\rm sym}^{\rm bulk}\delta^2.
\]
L'energia totale per barione e'
\[
\begin{aligned}
e&\simeq m_N-\frac{E_b}{A}\\
&\simeq m_N-a_v+a_{\rm sym}^{\rm bulk}\delta^2.
\end{aligned}
\]
Nel limite della materia nucleare: $a_v\simeq16$ MeV e $a_{\rm sym}^{\rm bulk}=J\simeq32$ MeV.
\medskip
\textit{Nota:} il coefficiente di simmetria di un nucleo finito non coincide esattamente con $J$ perche' include effetti di superficie.
\end{frame}

% 09
\begin{frame}{Energia per barione e potenziale chimico non coincidono sempre}
\small
Per materia fredda con un solo numero barionico conservato:
\[
\mu_B=\frac{d\varepsilon}{dn_B},\qquad P=n_B\mu_B-\varepsilon,
\]
quindi
\[
{\mu_B=\frac{\varepsilon}{n_B}+\frac{P}{n_B}=e+\frac{P}{n_B}.}
\]
\begin{columns}[T,onlytextwidth]
\begin{column}{0.49\textwidth}
Materia \emph{simmetrica} alla saturazione:
\[
P(n_{\rm sat},\delta=0)=0
\]
e percio' $\mu_B=e\simeq m_N-16$ MeV.
\end{column}
\begin{column}{0.49\textwidth}
Materia \emph{neutronica} alla stessa densita':
\[
P_{\rm PNM}(n_{\rm sat})>0
\]
in generale, dunque $\mu_n>e_{\rm PNM}$.
\end{column}
\end{columns}
\medskip
Il valore $955$ MeV del testo fornisce la \emph{scala} $E/A$, non necessariamente il potenziale chimico esatto.
\end{frame}

% 10
\begin{frame}{Correzione esplicita: il ruolo della pendenza $L$}
\small
Scriviamo
\[
e(n_B,\delta)\simeq e_{\rm sat}(n_B)+S(n_B)\delta^2.
\]
Per neutroni puri a $n_{\rm sat}$, $\delta=1$, mentre $e_{\rm sat}'(n_{\rm sat})=0$.
\[
\begin{aligned}
P_{\rm PNM}(n_{\rm sat})
&=n_{\rm sat}^2\left.\frac{d e_{\rm PNM}}{dn_B}\right|_{n_{\rm sat}}\\
&\simeq n_{\rm sat}^2 S'(n_{\rm sat}).
\end{aligned}
\]
Definendo $L=3n_{\rm sat}S'(n_{\rm sat})$:
\[
P_{\rm PNM}(n_{\rm sat})\simeq\frac{n_{\rm sat}L}{3},
\qquad
{\mu_n\simeq m_N+16+\frac L3.}
\]
La stima del libro $\mu_B\sim955$ MeV trascura questo termine di pressione.
\end{frame}

\subsection{Vincolo perturbativo ad altissime densita'}

% 11
\begin{frame}{Limite pQCD: quanto e' distante dalla saturazione?}
\small
\begin{columns}[T,onlytextwidth]
\begin{column}{0.49\textwidth}
Il testo prende come scala indicativa:
\[
\mu_q\simeq900\,\mathrm{MeV},
\]
con
\[
\mu_B=3\mu_q\simeq2.7\,\mathrm{GeV}.
\]
All'estremo nucleonico:
\[
\mu_B\sim0.95\,\mathrm{GeV}.
\]
Differenza nei potenziali chimici: meno di un fattore $3$.
\end{column}
\begin{column}{0.49\textwidth}
Ma la densita' di energia cresce approssimativamente come
\[
\varepsilon_{\rm free}\propto\mu_q^4.
\]
Ne consegue un divario assai maggiore fra densita' nucleari e densita' a cui la pQCD e' controllata.
\medskip
La materia nel core delle stelle osservate ricade in gran parte nella regione intermedia.
\end{column}
\end{columns}
\end{frame}

% 12
\begin{frame}{Densita' del gas libero a $\mu_q=900$ MeV}
\small
Tre sapori massless, $T=0$:
\[
P_{\rm SB}=\frac{3\mu_q^4}{4\pi^2},
\qquad \varepsilon_{\rm SB}=3P_{\rm SB}=\frac{9\mu_q^4}{4\pi^2}.
\]
Ripristinando $\hbar c$:
\[
\varepsilon_{\rm SB}[\mathrm{MeV\,fm^{-3}}]
=\frac{9(900\,\mathrm{MeV})^4}{4\pi^2(\hbar c)^3},
\quad \hbar c\simeq197.327\,\mathrm{MeV\,fm}.
\]
Numericalmente:
\[
\varepsilon_{\rm SB}\simeq1.95\times10^4\,\mathrm{MeV\,fm^{-3}}
\simeq140\,\varepsilon_0.
\]
Con le correzioni interagenti, il libro stima $\varepsilon_{\rm pQCD}\sim70\text{--}100\,\varepsilon_0$.
\end{frame}

\subsection{Figura 9.1: il divario fra le EOS}

% 13
\begin{frame}{Figura 9.1: EOS nucleonica, regione incerta e pQCD}
\small
\begin{columns}[T,onlytextwidth]
\begin{column}{0.63\textwidth}
\centering
\includegraphics[width=\linewidth]{hybrid9_fig91.png}
\end{column}
\begin{column}{0.35\textwidth}
\textbf{Asse verticale}
\[
P\ [\mathrm{MeV\,fm^{-3}}].
\]
\textbf{Asse orizzontale}
\[
\widetilde\mu_q=\mu_q-\frac{\mu_{\rm Fe}}3.
\]
Entrambi gli assi: scala logaritmica.
\medskip
A sinistra: croste e materia nucleonica. A destra: pQCD.
\end{column}
\end{columns}
\vspace{1mm}
{\scriptsize Figura 9.1 dal materiale allegato (Kurkela et al., 2014, come citato nel libro).}
\end{frame}

% 14
\begin{frame}{Perche' traslare il potenziale chimico dei quark?}
\small
Il potenziale chimico di riferimento del ferro e' assunto come
\[
\mu_{\rm Fe}\simeq930\,\mathrm{MeV},\qquad \frac{\mu_{\rm Fe}}3\simeq310\,\mathrm{MeV}.
\]
Definendo $\widetilde\mu_q=\mu_q-\mu_{\rm Fe}/3$,
\[
\mu_q=310\,\mathrm{MeV}+\widetilde\mu_q,
\qquad \mu_B=930\,\mathrm{MeV}+3\widetilde\mu_q.
\]
Esempio:
\[
\widetilde\mu_q=10\,\mathrm{MeV}
\quad\Longrightarrow\quad\mu_B\simeq960\,\mathrm{MeV}.
\]
La traslazione rende leggibili su scala logaritmica piccoli incrementi di potenziale chimico rispetto alla materia nucleare stabile.
\end{frame}

% 15
\begin{frame}{Lettura dei confini nella Figura 9.1}
\small
Le linee verticali delimitano indicativamente:
\[
\begin{array}{ll}
\widetilde\mu_q\simeq3\,\mathrm{MeV} & \text{fine crosta esterna},\\[2mm]
\widetilde\mu_q\simeq7.4\,\mathrm{MeV} & \text{fine crosta interna},\\[2mm]
\widetilde\mu_q\simeq16\,\mathrm{MeV} & \text{estremo nucleonico}.
\end{array}
\]
Il margine inferiore pQCD si trova molto piu' a destra ($\widetilde\mu_q\sim560$ MeV nel grafico).
\medskip
La zona centrale, con il punto interrogativo, richiede una EOS interpolata; non equivale a una regione sperimentalmente accertata di quark liberi.
\medskip
La fascia vicina all'asse orizzontale segnala che le densita' centrali delle stelle di massa massima ricadono proprio nell'intervallo incerto.
\end{frame}

\section{Interpolazione con politropiche a tratti}
\subsection{Termodinamica a temperatura zero}

% 16
\begin{frame}{Politropiche a tratti: definizione}
\small
Per ogni intervallo di densita' si assume
\[
{P(n)=Kn^{\Gamma}},
\]
con $n=n_B$, $K$ costante e $\Gamma$ indice politropico.
\[
\Gamma=\frac{d\ln P}{d\ln n}.
\]
In una EOS a tratti:
\[
P(n)=\begin{cases}
K_1n^{\Gamma_1},&n_0\leq n<n_1,\\
K_2n^{\Gamma_2},&n_1\leq n<n_2,\\
\cdots&\cdots
\end{cases}
\]
Le costanti di ciascun segmento devono essere determinate con condizioni di raccordo termodinamico.
\end{frame}

% 17
\begin{frame}{Identita' termodinamiche essenziali}
\small
A temperatura zero, materia neutra in equilibrio con un'unica carica conservata indipendente:
\[
d\varepsilon=\mu\,dn,\qquad
P=\mu n-\varepsilon,\qquad dP=n\,d\mu.
\]
Da esse seguono
\[
{\mu=\frac{d\varepsilon}{dn}},
\qquad
{n=\frac{dP}{d\mu}}.
\]
\begin{columns}[T,onlytextwidth]
\begin{column}{0.49\textwidth}
La forma $P(n)$ \emph{non basta} a definire autonomamente l'energia: occorre integrare l'identita' termodinamica.
\end{column}
\begin{column}{0.49\textwidth}
La continuita' di $P(\mu)$ e della sua derivata garantisce un raccordo senza salto della densita' barionica.
\end{column}
\end{columns}
\end{frame}

% 18
\begin{frame}{Derivazione dell'identita' (9.2)}
\small
Partiamo da
\[
P=n\mu-\varepsilon,
\qquad\mu=\frac{d\varepsilon}{dn}.
\]
Sostituiamo il potenziale chimico:
\[
P=n\frac{d\varepsilon}{dn}-\varepsilon.
\]
Dividendo per $n^2$:
\[
\frac P{n^2}=\frac1n\frac{d\varepsilon}{dn}-\frac{\varepsilon}{n^2}.
\]
Il membro destro e' la derivata di un rapporto:
\[
{P=n^2\frac{d}{dn}\left(\frac{\varepsilon}{n}\right).}
\]
\end{frame}

% 19
\begin{frame}{Dalla politropica alla densita' di energia (9.3)}
\small
Assumiamo in un intervallo
\[
P(n)=Kn^\Gamma.
\]
Dall'identita' precedente:
\[
Kn^\Gamma=n^2\frac{d}{dn}\left(\frac{\varepsilon}{n}\right).
\]
Dividiamo per $n^2$:
\[
{\frac{d}{dn}\left(\frac{\varepsilon}{n}\right)=Kn^{\Gamma-2}.}
\]
Per $\Gamma\ne1$, integriamo:
\[
\frac{\varepsilon}{n}=K\int n^{\Gamma-2}dn+c
=\frac{K}{\Gamma-1}n^{\Gamma-1}+c.
\]
\end{frame}

% 20
\begin{frame}{Densita' di energia politropica (9.4)}
\small
Moltiplicando il risultato integrato per $n$:
\[
{\varepsilon(n)=\frac{K}{\Gamma-1}n^\Gamma+cn.}
\]
Poiche' $P=Kn^\Gamma$,
\[
{\varepsilon(n)=\frac{P(n)}{\Gamma-1}+cn.}
\]
\begin{columns}[T,onlytextwidth]
\begin{column}{0.49\textwidth}
Il primo termine e' determinato dalla legge politropica $P(n)$.
\end{column}
\begin{column}{0.49\textwidth}
Il termine $cn$ integra l'informazione sull'energia per barione al raccordo; in generale $c\ne m_N$.
\end{column}
\end{columns}
\end{frame}

% 21
\begin{frame}{Potenziale chimico della politropica (9.5)}
\small
Per definizione:
\[
\mu(n)=\frac{d\varepsilon(n)}{dn}.
\]
Derivando termine a termine:
\[
\begin{aligned}
\mu(n)&=\frac{d}{dn}\left[\frac{K}{\Gamma-1}n^\Gamma+cn\right]\\
&=\frac{K\Gamma}{\Gamma-1}n^{\Gamma-1}+c.
\end{aligned}
\]
Quindi
\[
{\mu(n)=c+\frac{K\Gamma}{\Gamma-1}n^{\Gamma-1}.}
\]
Verifica: $n\mu-\varepsilon=Kn^\Gamma=P$.
\end{frame}

\subsection{Raccordo al punto $(n_0,P_0,\mu_0)$}

% 22
\begin{frame}{Primo vincolo di raccordo: pressione}
\small
Sia $(n_0,P_0,\mu_0)$ il punto terminale dell'EOS nota.
\medskip
Per il segmento politropico imponiamo
\[
P(n_0)=P_0.
\]
Con $P(n)=Kn^\Gamma$:
\[
P_0=Kn_0^\Gamma
\quad\Longrightarrow\quad
{K=\frac{P_0}{n_0^\Gamma}.}
\]
\begin{columns}[T,onlytextwidth]
\begin{column}{0.49\textwidth}
$\Gamma$ determina la forma del segmento.
\end{column}
\begin{column}{0.49\textwidth}
$K$ non e' libero una volta scelta la pressione al raccordo.
\end{column}
\end{columns}
\medskip
Il simbolo $n_0$ qui identifica il \emph{punto di raccordo}, non necessariamente la densita' di saturazione nucleare.
\end{frame}

% 23
\begin{frame}{Secondo vincolo: continuita' del potenziale (9.6)}
\small
Partiamo da
\[
\mu(n)=c+\frac{K\Gamma}{\Gamma-1}n^{\Gamma-1}.
\]
Imponiamo $\mu(n_0)=\mu_0$:
\[
\mu_0=c+\frac{K\Gamma}{\Gamma-1}n_0^{\Gamma-1}.
\]
Si ricava
\[
{c=\mu_0-\frac{K\Gamma}{\Gamma-1}n_0^{\Gamma-1}.}
\]
Inserendo $K=P_0/n_0^\Gamma$:
\[
{c=\mu_0-\frac{\Gamma}{\Gamma-1}\frac{P_0}{n_0}.}
\]
Le due condizioni fissano $K$ e $c$; rimane da scegliere $\Gamma$.
\end{frame}

% 24
\begin{frame}{Controllo del raccordo: anche $\varepsilon$ e' continua}
\small
Alla densita' $n_0$,
\[
\varepsilon(n_0)=\frac{P_0}{\Gamma-1}+cn_0.
\]
Usiamo
\[
cn_0=\mu_0n_0-\frac{\Gamma}{\Gamma-1}P_0.
\]
Segue
\[
\begin{aligned}
\varepsilon(n_0)
&=\frac{P_0}{\Gamma-1}+\mu_0n_0
-\frac{\Gamma P_0}{\Gamma-1}\\
&=\mu_0n_0-P_0.
\end{aligned}
\]
\[
{\varepsilon_0+P_0=\mu_0n_0.}
\]
La densita' di energia coincide con quella dell'EOS nota se il punto soddisfa la stessa identita' termodinamica.
\end{frame}

\subsection{Da $P(n)$ a $P(\mu)$ e $\varepsilon(P)$}

% 25
\begin{frame}{Invertire il potenziale chimico: primo passaggio}
\small
Partiamo dalla (9.5):
\[
\mu-c=\frac{K\Gamma}{\Gamma-1}n^{\Gamma-1}.
\]
Isolando la potenza di $n$:
\[
n^{\Gamma-1}=\frac{\Gamma-1}{K\Gamma}(\mu-c).
\]
Elevando a $1/(\Gamma-1)$:
\[
{n(\mu)=\left[\frac{\Gamma-1}{K\Gamma}(\mu-c)\right]^{1/(\Gamma-1)}.}
\]
L'argomento della potenza deve rimanere positivo nel tratto fisico scelto.
\end{frame}

% 26
\begin{frame}{Densita' in funzione di $\mu$: raccordo (9.7)}
\small
Sostituiamo
\[
c=\mu_0-\frac{K\Gamma}{\Gamma-1}n_0^{\Gamma-1}
\]
nel termine $\mu-c$:
\[
\mu-c=(\mu-\mu_0)+\frac{K\Gamma}{\Gamma-1}n_0^{\Gamma-1}.
\]
Pertanto
\[
{
 n(\mu)=\left[n_0^{\Gamma-1}
 +\frac{\Gamma-1}{K\Gamma}(\mu-\mu_0)
 \right]^{1/(\Gamma-1)}.}
\]
Controllo immediato: $n(\mu_0)=n_0$.
\end{frame}

% 27
\begin{frame}{Pressione in funzione di $\mu$: Eq. (9.8)}
\small
Partendo da $P=Kn^\Gamma$ e dalla (9.7):
\[
{
P(\mu)=K\left[
 n_0^{\Gamma-1}+\frac{\Gamma-1}{K\Gamma}(\mu-\mu_0)
\right]^{\Gamma/(\Gamma-1)}.}
\]
Al raccordo $\mu=\mu_0$:
\[
P(\mu_0)=Kn_0^\Gamma=P_0.
\]
La funzione pressione e' quindi continua. Rimane da verificarne la derivata.
\end{frame}

% 28
\begin{frame}{Verifica: $dP/d\mu=n$ in ogni segmento}
\small
Definiamo
\[
X(\mu)=n_0^{\Gamma-1}+\frac{\Gamma-1}{K\Gamma}(\mu-\mu_0).
\]
Allora $P(\mu)=KX^{\Gamma/(\Gamma-1)}$ e
\[
\begin{aligned}
\frac{dP}{d\mu}
&=K\frac{\Gamma}{\Gamma-1}X^{1/(\Gamma-1)}
\frac{\Gamma-1}{K\Gamma}\\
&=X^{1/(\Gamma-1)}=n(\mu).
\end{aligned}
\]
Per $\mu=\mu_0$:
\[
{\left.\frac{dP}{d\mu}\right|_{\mu_0}=n_0.}
\]
Questa e' la condizione sulla pendenza usata per il raccordo liscio delle politropiche.
\end{frame}

% 29
\begin{frame}{Dalla pressione alla densita' per le TOV}
\small
Per integrare le TOV occorre $\varepsilon(P)$.
\medskip
Da $P=Kn^\Gamma$:
\[
n=\left(\frac PK\right)^{1/\Gamma}.
\]
Poiche' $P_0=Kn_0^\Gamma$, equivalentemente
\[
{n(P)=n_0\left(\frac P{P_0}\right)^{1/\Gamma}.}
\]
Usiamo la densita' di energia gia' integrata:
\[
\varepsilon(n)=\frac{P(n)}{\Gamma-1}+cn.
\]
Sostituiremo $n(P)$ e poi il valore di $c$ fissato al raccordo.
\end{frame}

% 30
\begin{frame}{EOS per le TOV: derivazione di (9.9)}
\small
Sostituendo $n(P)$:
\[
\varepsilon(P)=\frac{P}{\Gamma-1}+cn_0\left(\frac{P}{P_0}\right)^{1/\Gamma}.
\]
Dal vincolo sul potenziale:
\[
cn_0=\mu_0n_0-\frac{\Gamma}{\Gamma-1}P_0.
\]
Si ottiene
\[
{
\begin{aligned}
\varepsilon(P)={}&\frac{P}{\Gamma-1}\\
&+\left(\mu_0n_0-\frac{\Gamma}{\Gamma-1}P_0\right)
\left(\frac{P}{P_0}\right)^{1/\Gamma}.
\end{aligned}}
\]
Il termine in $P^{1/\Gamma}$ conserva l'informazione termodinamica del raccordo.
\end{frame}

% 31
\begin{frame}{Il caso speciale $\Gamma=1$}
\small
Le formule precedenti con $1/(\Gamma-1)$ richiedono $\Gamma\ne1$.
\medskip
Se $P=Kn$, l'equazione (9.2) diventa
\[
\frac{d}{dn}\left(\frac{\varepsilon}{n}\right)=\frac Kn.
\]
Integrando, con densita' di riferimento $n_{\rm ref}$,
\[
{\varepsilon(n)=Kn\ln\frac{n}{n_{\rm ref}}+cn.}
\]
Derivando:
\[
{\mu(n)=K\left[\ln\frac{n}{n_{\rm ref}}+1\right]+c.}
\]
Il segmento con $\Gamma=1$ si tratta dunque separatamente mediante un logaritmo.
\end{frame}

% 32
\begin{frame}{Algoritmo di raccordo per piu' politropiche}
\small
\begin{enumerate}
\item Partire dal punto nucleonico $(n_0,P_0,\mu_0)$.
\item Scegliere $\Gamma_1$ e determinare $K_1,c_1$ imponendo $P_0,\mu_0$.
\item Proseguire fino a un nuovo punto $(n_1,P_1,\mu_1)$.
\item Scegliere $\Gamma_2$ e imporre le stesse condizioni per il secondo segmento.
\item Raccordare all'EOS pQCD, controllando vincoli fisici lungo tutta la curva.
\end{enumerate}
\medskip
Al raccordo liscio:
\[
P_-(\mu_i)=P_+(\mu_i),\qquad
\left.\frac{dP_-}{d\mu}\right|_{\mu_i}
=\left.\frac{dP_+}{d\mu}\right|_{\mu_i}=n_i.
\]
\end{frame}

\subsection{Causalita', rigidita' e transizioni di fase}

% 33
\begin{frame}{Derivazione della velocita' del suono}
\small
A $T=0$, in un fluido barotropico:
\[
c_s^2=\frac{dP}{d\varepsilon}
=\frac{dP/dn}{d\varepsilon/dn}
\qquad(c=1).
\]
Per $P=Kn^\Gamma$:
\[
\frac{dP}{dn}=\Gamma Kn^{\Gamma-1}=\frac{\Gamma P}{n}.
\]
Dalla termodinamica:
\[
\frac{d\varepsilon}{dn}=\mu=\frac{\varepsilon+P}{n}.
\]
Quindi
\[
{c_s^2=\frac{\Gamma P}{\varepsilon+P}.}
\]
Causalita' e stabilita' locale richiedono $0\leq c_s^2\leq1$.
\end{frame}

% 34
\begin{frame}{Raccordo liscio e transizione del primo ordine}
\small
\begin{columns}[T,onlytextwidth]
\begin{column}{0.49\textwidth}
\textbf{Interpolazione liscia}
\[
P_H(\mu_i)=P_Q(\mu_i),
\]
\[
n_H(\mu_i)=n_Q(\mu_i).
\]
Per $\varepsilon=\mu n-P$ risulta anche $\varepsilon_H=\varepsilon_Q$ al punto.
\medskip
Nessun salto di densita' introdotto dal raccordo.
\end{column}
\begin{column}{0.49\textwidth}
\textbf{Primo ordine (Maxwell)}
\[
P_H(\mu_t)=P_Q(\mu_t),
\]
\[
n_Q\ne n_H.
\]
A pressione e potenziale comuni,
\[
{\Delta\varepsilon=\mu_t\Delta n_B.}
\]
La pressione continua, ma la derivata $dP/d\mu$ salta.
\end{column}
\end{columns}
\end{frame}

% 35
\begin{frame}{Quali politropiche sono fisicamente ammissibili?}
\small
\begin{columns}[T,onlytextwidth]
\begin{column}{0.49\textwidth}
\textbf{Vincoli locali}
\begin{itemize}
\item $n=dP/d\mu>0$;
\item stabilita' meccanica $dP/dn>0$;
\item causalita' $c_s^2=dP/d\varepsilon\leq1$.
\end{itemize}
\end{column}
\begin{column}{0.49\textwidth}
\textbf{Vincolo globale}
\[
M_{\max}\gtrsim2M_\odot
\]
per essere compatibili con l'esistenza di stelle massicce osservate.
\medskip
Una EOS troppo soffice non sostiene grandi masse; una troppo rigida puo' violare la causalita'.
\end{column}
\end{columns}
\medskip
La rigidita' stellare si giudica da $P(\varepsilon)$ e dalla velocita' del suono, non dal solo valore di $\Gamma$ isolatamente.
\end{frame}

\section{Massa--raggio delle stelle ibride}
\subsection{Figura 9.2 e vincoli astrofisici}

% 36
\begin{frame}{Figura 9.2: banda delle relazioni massa--raggio}
\small
\begin{columns}[T,onlytextwidth]
\begin{column}{0.64\textwidth}
\centering
\includegraphics[width=\linewidth]{hybrid9_fig92.png}
\end{column}
\begin{column}{0.34\textwidth}
Banda grigia: famiglie di EOS interpolate fra limite nucleonico e pQCD.
\medskip
Curve I, II e III: tre politropiche rappresentative.
\medskip
Linee: buchi neri, Buchdahl, limite causale e compattezze costanti.
\end{column}
\end{columns}
\vspace{1mm}
{\scriptsize Figura 9.2 dal materiale allegato (Kurkela et al., 2014, come citato nel libro).}
\end{frame}

% 37
\begin{frame}{Tre sequenze rappresentative nella Figura 9.2}
\footnotesize
\begin{columns}[T,onlytextwidth]
\begin{column}{0.33\textwidth}
\textbf{I: piu' compatta}
\medskip
Raggi relativamente piccoli.
\[
\varepsilon_{c,\max}\sim10\varepsilon_0.
\]
La configurazione di massa massima raggiunge densita' centrali elevate.
\end{column}
\begin{column}{0.33\textwidth}
\textbf{II: massa massima maggiore}
\[
M_{\max}\simeq2.4M_\odot,
\]
\[
\varepsilon_{c,\max}\sim5\varepsilon_0.
\]
Mostra come una EOS adeguatamente rigida possa sostenere masse elevate.
\end{column}
\begin{column}{0.33\textwidth}
\textbf{III: raggi maggiori}
\medskip
Sequenza piu' a destra nel diagramma.
\[
\varepsilon_{c,\max}\sim5\varepsilon_0.
\]
Raggi piu' ampi alle masse considerate.
\end{column}
\end{columns}
\end{frame}

% 38
\begin{frame}{Densita' centrali: cosa sondano le osservazioni?}
\small
Per la famiglia di EOS del testo:
\[
\begin{array}{c|c|c}
M & R & \varepsilon_c/\varepsilon_0\\ \hline
1.4M_\odot & 11\text{--}14.5\,\mathrm{km} & 1.6\text{--}4\\[1mm]
2.0M_\odot & 10\text{--}15\,\mathrm{km} & 2\text{--}8
\end{array}
\]
Confronto fra scale:
\[
\varepsilon_c\sim\text{pochi }\varepsilon_0,
\qquad \varepsilon_{\rm pQCD}\gtrsim70\,\varepsilon_0.
\]
Le osservazioni delle stelle compatte mettono alla prova la regione intermedia dell'EOS, non direttamente il regime di pQCD asintotica.
\medskip
{\scriptsize Intervalli del lavoro citato nel libro; non sono vincoli osservativi aggiornati.}
\end{frame}

% 39
\begin{frame}{Compattezza, causalita' e massa osservata}
\small
\begin{columns}[T,onlytextwidth]
\begin{column}{0.49\textwidth}
\textbf{Compattezza}
\[
\mathcal C=\frac{GM}{Rc^2}.
\]
Rette di riferimento nel grafico:
\[
\mathcal C=0.1,\ 0.2,\ 0.3.
\]
Limiti esaminati nel capitolo 8:
\[
\mathcal C_{\rm BH}=\tfrac12,\quad
\mathcal C_{\rm Buchdahl}=\tfrac49.
\]
\end{column}
\begin{column}{0.49\textwidth}
\textbf{Vincoli sull'EOS}
\[
c_s^2\leq1,\qquad M_{\max}\geq2M_\odot.
\]
La causalita' impedisce una EOS arbitrariamente rigida.
\medskip
La massa massima osservata impedisce una EOS arbitrariamente soffice.
\medskip
La loro combinazione restringe la banda ammissibile nel diagramma $M$--$R$.
\end{column}
\end{columns}
\end{frame}

\subsection{Riepilogo}

% 40
\begin{frame}{Risultato del \S9.1: interpolare senza inventare la microfisica}
\small
\begin{enumerate}
\item EOS nucleonica vincolata vicino a $\varepsilon_0$.
\item EOS pQCD controllata solo intorno a $\varepsilon\gtrsim70\text{--}100\varepsilon_0$.
\item Politropiche a tratti raccordate con continuita' di $P$, $\mu$ e $n$.
\item Verifica di causalita', stabilita' e massa massima.
\item Integrazione delle TOV per ottenere una banda di curve $M(R)$.
\end{enumerate}
\medskip
Il raccordo liscio e' una \emph{parametrizzazione} dell'EOS intermedia, non una dimostrazione che la transizione adroni--quark sia un crossover.
\end{frame}

% 41
\begin{frame}{Formule da ricordare: politropiche termodinamiche}
\footnotesize
\begin{columns}[T,onlytextwidth]
\begin{column}{0.49\textwidth}
\[
P(n)=Kn^\Gamma,
\]
\[
\varepsilon(n)=\frac{Kn^\Gamma}{\Gamma-1}+cn,
\]
\[
\mu(n)=\frac{K\Gamma}{\Gamma-1}n^{\Gamma-1}+c,
\]
\[
c=\mu_0-\frac{\Gamma P_0}{(\Gamma-1)n_0}.
\]
\end{column}
\begin{column}{0.49\textwidth}
\[
n=\frac{dP}{d\mu},
\]
\[
\varepsilon(P)=\frac{P}{\Gamma-1}
+\left(\mu_0n_0-\frac{\Gamma P_0}{\Gamma-1}\right)
\left(\frac{P}{P_0}\right)^{1/\Gamma},
\]
\[
c_s^2=\frac{\Gamma P}{\varepsilon+P}\leq1.
\]
\end{column}
\end{columns}
\medskip
\textbf{Energia di simmetria:} $e_{\rm PNM}(n_{\rm sat})\simeq m_N-16+32=m_N+16\simeq955$ MeV.
\end{frame}

\section{Transizioni di fase nella materia densa}

% 01
\begin{frame}{Dal raccordo liscio alle transizioni di fase}
\small
Nel \S 9.1 l'EOS intermedia era ottenuta interpolando fra
materia nucleonica e materia di quark con raccordi lisci.
\medskip
Nel \S 9.2 ammettiamo una transizione di fase del primo ordine:
\[
\text{fase adronica I}\quad\longleftrightarrow\quad
\text{fase di quark II}.
\]
\begin{columns}[T,onlytextwidth]
\begin{column}{0.49\textwidth}
\textbf{Domande termodinamiche}
\begin{itemize}
\item Quando coesistono le due fasi?
\item Quali grandezze sono continue?
\item E' consentito un salto di densita'?
\end{itemize}
\end{column}
\begin{column}{0.49\textwidth}
\textbf{Conseguenze stellari}
\begin{itemize}
\item Interfaccia netta oppure fase mista estesa.
\item Possibili domini caricati (``pasta'').
\item Modifiche all'EOS e alle soluzioni TOV.
\end{itemize}
\end{column}
\end{columns}
\end{frame}

\subsection{Diagramma di fase della QCD}

% 02
\begin{frame}{Diagramma di fase della QCD: figura 9.3}
\small
\begin{columns}[T,onlytextwidth]
\begin{column}{0.64\textwidth}
    \begin{figure}[!ht]\includegraphics[trim={0cm 0cm 0 0},clip,keepaspectratio,width=0.95\textwidth]{Schaffner/QCD_phasediagram_T_mu}\label{fig:Schaffner_QCD_phasediagram_T_mu}\end{figure}
\end{column}
\begin{column}{0.34\textwidth}
Assi: temperatura $T$ e potenziale chimico di quark $\mu_q$.
\medskip
\begin{itemize}
\item Adroni: bassi $T$ e $\mu_q$.
\item Plasma di quark e gluoni: $T$ elevata.
\item Possibile pairing di quark: bassa $T$, elevata $\mu_q$.
\end{itemize}
\end{column}
\end{columns}
\medskip
La topologia delle linee di transizione ad alto $\mu_q$ e' \emph{congetturale}.
\end{frame}

% 03
\begin{frame}{Fase adronica, fase deconfinata e pairing}
\small
\begin{columns}[T,onlytextwidth]
\begin{column}{0.49\textwidth}
\textbf{Regime adronico}
\begin{itemize}
\item Quark confinati negli adroni.
\item A basse densita' e basse $T$, simmetria chirale spontaneamente rotta.
\end{itemize}
\[
\langle\bar q q\rangle\ne0.
\]
\textbf{Regime caldo}
\begin{itemize}
\item Quark e gluoni sono i gradi di liberta' rilevanti.
\item Il condensato chirale si riduce.
\end{itemize}
\end{column}
\begin{column}{0.49\textwidth}
\textbf{Regime freddo e denso}
\begin{itemize}
\item Superfici di Fermi dei quark.
\item Possibile condensazione di coppie $qq$.
\end{itemize}
\[
\langle qq\rangle\ne0
\quad\text{(simbolicamente)}.
\]
Il pairing di colore e' analogo al meccanismo BCS; esempi di fasi: 2SC e CFL.
\medskip
Un condensato $qq$ va definito con gli opportuni indici di colore, sapore e spin.
\end{column}
\end{columns}
\end{frame}

% 04
\begin{frame}{Rottura chirale e confinamento: concetti distinti}
\small
\begin{columns}[T,onlytextwidth]
\begin{column}{0.49\textwidth}
\textbf{Simmetria chirale}
\[
SU(N_f)_L\times SU(N_f)_R.
\]
Un indicatore della rottura spontanea e'
\[
\langle\bar q q\rangle\ne0.
\]
Riguarda il comportamento dei campi sinistri e destri rispetto a rotazioni indipendenti in sapore.
\end{column}
\begin{column}{0.49\textwidth}
\textbf{Confinamento}
\medskip
Riguarda l'assenza di particelle isolate con carica di colore tra gli stati asintotici.
\medskip
Il loop di Polyakov e' un parametro d'ordine rigoroso per la teoria di gauge pura, non per QCD con quark dinamici fisici.
\end{column}
\end{columns}
\medskip
Il ripristino chirale e il deconfinamento possono essere correlati,
ma non sono necessariamente un'unica transizione.
\end{frame}

% 05
\begin{frame}{Crossover, transizione del primo ordine, punto critico}
\small
\begin{columns}[T,onlytextwidth]
\begin{column}{0.49\textwidth}
\textbf{Crossover}
\begin{itemize}
\item Cambiamento rapido, ma continuo, delle osservabili termodinamiche.
\item A $\mu_B\simeq0$ e masse fisiche dei quark, il passaggio termico della QCD e' un crossover.
\end{itemize}
\[
T_c\sim150\,\mathrm{MeV}.
\]
\end{column}
\begin{column}{0.49\textwidth}
\textbf{Primo ordine}
\begin{itemize}
\item Coesistenza di due fasi.
\item Discontinuita' di densita' o di una derivata prima del potenziale termodinamico.
\item Una possibile linea puo' terminare in un punto critico.
\end{itemize}
\end{column}
\end{columns}
\medskip
La presenza del punto critico QCD e di un punto triplo nel disegno
non e' una previsione sperimentalmente stabilita.
\end{frame}

% 06
\begin{frame}{Quali regioni del diagramma si possono esplorare?}
\small
\begin{columns}[T,onlytextwidth]
\begin{column}{0.49\textwidth}
\textbf{Sistemi fisici}
\begin{itemize}
\item Universo primordiale: $T$ alta, $\mu_B$ modesto.
\item Collisioni di ioni pesanti: $T$ alta e $\mu_B$ variabile.
\item Supernove e merger: condizioni calde e dense.
\item Stelle di neutroni mature: $T$ bassa e $\mu_B$ elevato.
\end{itemize}
\end{column}
\begin{column}{0.49\textwidth}
\textbf{Distinguere due transizioni}
\medskip
La transizione liquido--gas \emph{nucleare} ha un punto critico stimato intorno a
\[
T_c^{\rm LG}\simeq17.9\,\mathrm{MeV},
\quad n_c\simeq0.06\,\mathrm{fm}^{-3}.
\]
Non coincide con l'eventuale transizione adroni--quark ad alta densita'.
\end{column}
\end{columns}
\end{frame}

\subsection{Costruzione di Maxwell}

% 07
\begin{frame}{Costruzione di Maxwell: ipotesi}
\small
Consideriamo due fasi omogenee a temperatura zero:
\[
\underbrace{\mathrm{I}}_{\text{adronica}}
\quad\text{e}\quad
\underbrace{\mathrm{II}}_{\text{di quark}}.
\]
Nel limite Maxwell si impone la \emph{neutralita' elettrica locale}
di ciascuna fase, riducendo il problema bulk a un potenziale
chimico indipendente, $\mu_B$.
\medskip
Cerchiamo le condizioni che rendono possibile la coesistenza
fra i due mezzi ai lati dell'interfaccia.
\begin{itemize}
\item Equilibrio meccanico: nessuna forza netta sull'interfaccia.
\item Equilibrio chimico: nessun trasferimento netto di numero barionico.
\end{itemize}
\end{frame}

% 08
\begin{frame}{Equilibrio meccanico e chimico}
\small
\begin{columns}[T,onlytextwidth]
\begin{column}{0.49\textwidth}
\textbf{Equilibrio meccanico}
\[
P_{\mathrm I}=P_{\mathrm{II}}.
\]
Se $P_{\mathrm I}>P_{\mathrm{II}}$, la prima fase tende a espandersi a danno della seconda.
\medskip
L'uguaglianza e' esatta per un'interfaccia bulk piana trascurandone la curvatura.
\end{column}
\begin{column}{0.49\textwidth}
\textbf{Equilibrio chimico}
\[
\mu_B^{\mathrm I}=\mu_B^{\mathrm{II}}.
\]
Se un barione passa da I a II, a $T,V$ fissati la variazione del potenziale e'
\[
\delta\Omega
=(\mu_B^{\mathrm{II}}-\mu_B^{\mathrm I})\,\delta N_B.
\]
All'equilibrio il trasferimento non favorisce alcuna direzione.
\end{column}
\end{columns}
\end{frame}

% 09
\begin{frame}{Condizione di coesistenza: equazione (9.10)}
\small
Introduciamo il potenziale chimico della transizione $\mu_t$.
Le pressioni delle due fasi, calcolate allo stesso $\mu_t$, sono uguali:
\[
P_{\mathrm I}(\mu_t)=P_{\mathrm{II}}(\mu_t),
\]
\[
\mu_B^{\mathrm I}=\mu_B^{\mathrm{II}}=\mu_t.
\]
\begin{columns}[T,onlytextwidth]
\begin{column}{0.49\textwidth}
Il punto di coesistenza si trova \emph{intersecando} le due curve $P_{\mathrm I}(\mu_B)$ e $P_{\mathrm{II}}(\mu_B)$.
\end{column}
\begin{column}{0.49\textwidth}
La costruzione non richiede
\[
n_B^{\mathrm I}=n_B^{\mathrm{II}},
\]
quindi puo' esserci un salto di densita'.
\end{column}
\end{columns}
\end{frame}

% 10
\begin{frame}{Figura 9.4: due rappresentazioni della stessa transizione}
\begin{figure}[!ht]\includegraphics[trim={0cm 0cm 0 0},clip,keepaspectratio,width=0.95\textwidth]{Schaffner/firsttransition_P_epsilon-mu}\label{fig:Schaffner_firsttransition_P_epsilon-mu}\end{figure}
\small
\begin{columns}[T,onlytextwidth]
\begin{column}{0.49\textwidth}
(a) $P(\varepsilon)$: la linea orizzontale A--D sostituisce la regione instabile/metastabile.
\end{column}
\begin{column}{0.49\textwidth}
(b) $P(\mu_B)$: al punto di incrocio le pressioni coincidono, ma non necessariamente le pendenze.
\end{column}
\end{columns}
\end{frame}

% 11
\begin{frame}{Quale fase e' termodinamicamente favorita?}
\small
Il potenziale gran-canonico a $T=0$ e'
\[
\Omega=E-\mu_B N_B,
\qquad \Omega=-PV
\quad\text{(bulk omogeneo)}.
\]
A $T,\mu_B,V$ fissati, l'equilibrio minimizza $\Omega$:
\[
\Omega_{\mathrm I}<\Omega_{\mathrm{II}}
\quad\Longleftrightarrow\quad
P_{\mathrm I}>P_{\mathrm{II}}.
\]
Quindi, alla stessa $\mu_B$, \emph{la fase a pressione maggiore e' favorita}.
\medskip
\begin{itemize}
\item Prima dell'intersezione: prevale il ramo stabile della fase I.
\item Dopo l'intersezione: prevale il ramo stabile della fase II.
\item All'intersezione: le due fasi coesistono.
\end{itemize}
\end{frame}

% 12
\begin{frame}{Stabilita' meccanica e ramo a pendenza negativa}
\small
Per un sistema a un numero conservato indipendente:
\[
d\varepsilon=\mu_B\,dn_B,
\qquad dP=n_B\,d\mu_B.
\]
Quindi
\[
\frac{dP}{d\varepsilon}
=\frac{n_B}{\mu_B}\frac{d\mu_B}{dn_B}
\equiv c_s^2
\quad(c=1).
\]
Nel ramo B--C della figura 9.4(a):
\[
\frac{dP}{d\varepsilon}<0
\quad\Longrightarrow\quad
\frac{d\mu_B}{dn_B}<0.
\]
La compressibilita' del ramo omogeneo e' negativa: instabilita' meccanica.
\medskip
La costruzione elimina anche tratti \emph{metastabili}, non solo il segmento B--C.
\end{frame}

% 13
\begin{frame}{Salto di densita' al potenziale di coesistenza}
\small
Per una sola carica conservata:
\[
n_B=\frac{dP}{d\mu_B}.
\]
Al punto di coesistenza, le due pressioni coincidono:
\[
P_{\mathrm I}(\mu_t)=P_{\mathrm{II}}(\mu_t)=P_t,
\]
ma le pendenze possono essere differenti:
\[
\left.\frac{dP_{\mathrm I}}{d\mu_B}\right|_{\mu_t}
=n_{\mathrm I}\ne n_{\mathrm{II}}
=\left.\frac{dP_{\mathrm{II}}}{d\mu_B}\right|_{\mu_t}.
\]
Il salto $\Delta n_B=n_{\mathrm{II}}-n_{\mathrm I}$ distingue
la transizione di primo ordine.
\end{frame}

% 14
\begin{frame}{Derivazione del salto di energia}
\small
A $T=0$, l'identita' di Gibbs--Duhem porta a
\[
\varepsilon+P=\mu_B n_B.
\]
Alle due estremita' del plateau:
\[
\varepsilon_{\mathrm I}=\mu_t n_{\mathrm I}-P_t,
\qquad
\varepsilon_{\mathrm{II}}=\mu_t n_{\mathrm{II}}-P_t.
\]
Sottraendo membro a membro:
\[
\begin{aligned}
\Delta\varepsilon
&=\varepsilon_{\mathrm{II}}-\varepsilon_{\mathrm I}\\
&=(\mu_t n_{\mathrm{II}}-P_t)
 -(\mu_t n_{\mathrm I}-P_t)\\
&=\mu_t\,(n_{\mathrm{II}}-n_{\mathrm I}).
\end{aligned}
\]
\[
\Delta\varepsilon=\mu_t\Delta n_B.
\]
\end{frame}

% 15
\begin{frame}{Il plateau di Maxwell nel piano $P$--$\varepsilon$}
\small
Sia $\chi$ la frazione volumetrica della fase II.
Trascurando le energie finite di superficie e Coulomb:
\[
\varepsilon_{\rm mix}=(1-\chi)\varepsilon_{\mathrm I}
+\chi\varepsilon_{\mathrm{II}},\qquad 0\le\chi\le1.
\]
Durante la coesistenza, nel limite bulk Maxwell,
\[
P_{\rm mix}=P_t,\qquad \mu_B=\mu_t.
\]
\begin{columns}[T,onlytextwidth]
\begin{column}{0.49\textwidth}
Variando $\chi$ fra $0$ e $1$, la densita' di energia media passa da $\varepsilon_{\mathrm I}$ a $\varepsilon_{\mathrm{II}}$.
\end{column}
\begin{column}{0.49\textwidth}
La pressione rimane costante: il ramo termodinamico e' orizzontale in $P(\varepsilon)$.
\end{column}
\end{columns}
\end{frame}

% 16
\begin{frame}{Perche' Maxwell produce un'interfaccia stellare netta?}
\small
Per una configurazione TOV statica, $P(r)$ diminuisce verso l'esterno:
\[
\frac{dP}{dr}
=-\frac{G(\varepsilon+P)
(m+4\pi r^3P/c^2)}
{c^2r^2[1-2Gm/(rc^2)]}<0.
\]
Uno strato macroscopico sul plateau Maxwell avrebbe invece
\[
P(r)=P_t\quad\Longrightarrow\quad \frac{dP}{dr}=0.
\]
La soluzione ideale passa quindi, a un raggio $r_t$, fra
\[
P(r_t^-)=P(r_t^+)=P_t,
\quad\varepsilon(r_t^-)\ne\varepsilon(r_t^+).
\]
Pressione continua, densita' di energia discontinua;
lo spessore microscopico dell'interfaccia e' qui trascurato.
\end{frame}

\subsection{Costruzione di Gibbs}

% 17
\begin{frame}{Due cariche conservate nella materia stellare}
\small
La materia fredda in equilibrio debole conserva almeno:
\[
N_B\quad\text{(numero barionico)},
\qquad N_Q\quad\text{(carica elettrica)}.
\]
Introduciamo i potenziali chimici coniugati
\[
\mu_B,\qquad\mu_Q.
\]
Per ogni specie $i$ in equilibrio:
\[
\mu_i=B_i\mu_B+Q_i\mu_Q.
\]
\begin{columns}[T,onlytextwidth]
\begin{column}{0.49\textwidth}
$B_i$: numero barionico della specie.
\end{column}
\begin{column}{0.49\textwidth}
$Q_i$: carica elettrica in unita' di $e$.
\end{column}
\end{columns}
\medskip
Ora la pressione di ciascuna fase e', in generale,
$P(\mu_B,\mu_Q)$.
\end{frame}

% 18
\begin{frame}{Esempio: neutroni, protoni ed elettroni}
\small
Per $B_n=B_p=1$, $B_e=0$ e $Q_n=0$, $Q_p=+1$, $Q_e=-1$:
\[
\begin{aligned}
\mu_n&=\mu_B,\\
\mu_p&=\mu_B+\mu_Q,\\
\mu_e&=-\mu_Q.
\end{aligned}
\]
L'equilibrio beta per $n\leftrightarrow p+e^-+\bar\nu_e$
(con neutrini liberi, $\mu_{\nu_e}=0$) segue automaticamente:
\[
\begin{aligned}
\mu_p+\mu_e
&=(\mu_B+\mu_Q)-\mu_Q\\
&=\mu_B=\mu_n.
\end{aligned}
\]
Le cariche conservate determinano quindi le relazioni fra i potenziali chimici.
\end{frame}

% 19
\begin{frame}{Differenziale della pressione con due cariche}
\small
Per la densita' di energia $\varepsilon(n_B,n_Q)$, a $T=0$:
\[
d\varepsilon=\mu_B\,dn_B+\mu_Q\,dn_Q.
\]
La trasformata di Legendre fornisce
\[
P=-\varepsilon+\mu_Bn_B+\mu_Qn_Q.
\]
Differenziando:
\[
\begin{aligned}
dP&=-d\varepsilon+n_Bd\mu_B+\mu_Bdn_B\\
  &\quad+n_Qd\mu_Q+\mu_Qdn_Q\\
&=n_Bd\mu_B+n_Qd\mu_Q.
\end{aligned}
\]
Quindi
\[
n_B=\frac{\partial P}{\partial\mu_B},
\qquad n_Q=\frac{\partial P}{\partial\mu_Q}.
\]
\end{frame}

% 20
\begin{frame}{Condizioni di Gibbs: equazione (9.11)}
\small
Per fasi I e II in equilibrio bulk a temperatura nulla:
\[
\begin{gathered}
P_{\mathrm I}(\mu_B,\mu_Q)
=P_{\mathrm{II}}(\mu_B,\mu_Q),\\[5pt]
\mu_B^{\mathrm I}=\mu_B^{\mathrm{II}},
\qquad
\mu_Q^{\mathrm I}=\mu_Q^{\mathrm{II}}.
\end{gathered}
\]
\begin{itemize}
\item Un solo valore di $\mu_B$ e $\mu_Q$ determina entrambe le composizioni.
\item Le due fasi possono avere cariche elettriche \emph{non nulle e opposte}.
\item La neutralita' e' imposta sulla miscela complessiva.
\end{itemize}
\medskip
\footnotesize
Per domini caricati, l'equilibrio elettrico va formulato rigorosamente
con i potenziali \emph{elettrochimici}; qui usiamo il limite bulk del testo.
\end{frame}

% 21
\begin{frame}{Figura 9.5: due superfici di pressione}
\small
\begin{columns}[T,onlytextwidth]
\begin{column}{0.60\textwidth}
    \begin{figure}[!ht]\includegraphics[trim={0cm 0cm 0 0},clip,keepaspectratio,width=0.95\textwidth]{Schaffner/Gibbs_nuclear_quark}\label{fig:Schaffner_Gibbs_nuclear_quark}\end{figure}
\end{column}
\begin{column}{0.38\textwidth}
\begin{itemize}
\item HP: fase adronica.
\item QP: fase di quark.
\item MP: \emph{linea} di coesistenza delle pressioni.
\end{itemize}
\medskip
Le curve HP e QP individuano la neutralita' della singola fase.
\end{column}
\end{columns}
\end{frame}

% 22
\begin{frame}{Perche' la coesistenza di Gibbs e' una linea?}
\small
\textbf{Maxwell:} dopo la neutralita' locale resta una variabile,
$\mu_B$:
\[
P_{\mathrm I}(\mu_B)=P_{\mathrm{II}}(\mu_B).
\]
L'equazione individua, normalmente, un punto $\mu_t$.
\medskip
\textbf{Gibbs:} le pressioni dipendono da due variabili:
\[
P_{\mathrm I}(\mu_B,\mu_Q)
=P_{\mathrm{II}}(\mu_B,\mu_Q).
\]
Una condizione per due potenziali chimici identifica in generale
una curva nel piano $(\mu_B,\mu_Q)$.
\medskip
In un grafico tridimensionale $(\mu_B,\mu_Q,P)$, la curva
corrisponde all'intersezione di due superfici.
\end{frame}

% 23
\begin{frame}{Neutralita' globale: equazione (9.12)}
\small
Sia $\chi$ la frazione di volume della fase II:
\[
\chi=\frac{V_{\mathrm{II}}}{V_{\mathrm I}+V_{\mathrm{II}}},
\qquad 0\le\chi\le1.
\]
La densita' di carica totale e' la media volumetrica:
\[
n_Q^{\rm tot}=(1-\chi)n_Q^{\mathrm I}
+\chi n_Q^{\mathrm{II}}.
\]
La neutralita' globale richiede
\[
(1-\chi)n_Q^{\mathrm I}+\chi n_Q^{\mathrm{II}}=0.
\]
I valori $n_Q^{\mathrm I}$ e $n_Q^{\mathrm{II}}$ vanno calcolati
allo \emph{stesso} $\mu_B,\mu_Q$ lungo la linea di coesistenza.
\medskip
\footnotesize
Se non sono inclusi nelle singole densita' di fase, i leptoni
contribuiscono con un termine separato $n_Q^{\ell}$.
\end{frame}

% 24
\begin{frame}{Derivazione della frazione volumetrica}
\small
Partiamo dalla neutralita' globale:
\[
(1-\chi)n_Q^{\mathrm I}+\chi n_Q^{\mathrm{II}}=0.
\]
Sviluppiamo e raccogliamo $\chi$:
\[
\begin{aligned}
n_Q^{\mathrm I}-\chi n_Q^{\mathrm I}
+\chi n_Q^{\mathrm{II}}&=0,\\
\chi(n_Q^{\mathrm{II}}-n_Q^{\mathrm I})&=-n_Q^{\mathrm I}.
\end{aligned}
\]
Quindi
\[
\chi=\frac{n_Q^{\mathrm I}}
{n_Q^{\mathrm I}-n_Q^{\mathrm{II}}}.
\]
Per avere $0\le\chi\le1$, le due cariche devono potersi compensare.
\end{frame}

% 25
\begin{frame}{Esempio numerico di neutralita' globale}
\small
Supponiamo, alla stessa coppia $(\mu_B,\mu_Q)$:
\[
n_Q^{\mathrm I}=+0.02\,\mathrm{fm}^{-3},
\qquad n_Q^{\mathrm{II}}=-0.03\,\mathrm{fm}^{-3}.
\]
La frazione della fase di quark necessaria alla neutralita' e'
\[
\begin{aligned}
\chi&=\frac{n_Q^{\mathrm I}}{n_Q^{\mathrm I}-n_Q^{\mathrm{II}}}\\
&=\frac{0.02}{0.02+0.03}=0.4.
\end{aligned}
\]
Pertanto
\[
V_{\mathrm{II}}=0.4V_{\rm tot},
\qquad V_{\mathrm I}=0.6V_{\rm tot}.
\]
Le due regioni possono essere cariche, pur essendo
neutro il campione complessivo.
\end{frame}

% 26
\begin{frame}{Energia e numero barionico: equazione (9.13)}
\small
Nel limite bulk, prima di introdurre tensione superficiale
e costo elettrostatico, le densita' medie sono
\[
\varepsilon_{\rm tot}=(1-\chi)\varepsilon_{\mathrm I}
+\chi\varepsilon_{\mathrm{II}},
\]
\[
n_B^{\rm tot}=(1-\chi)n_B^{\mathrm I}
+\chi n_B^{\mathrm{II}}.
\]
\begin{columns}[T,onlytextwidth]
\begin{column}{0.49\textwidth}
Al punto $1$ della figura 9.5:
\[
\chi\simeq0
\]
prevale la fase adronica.
\end{column}
\begin{column}{0.49\textwidth}
Al punto $2$ della figura 9.5:
\[
\chi\simeq1
\]
prevale la fase di quark.
\end{column}
\end{columns}
\medskip
Nella fase mista $\chi$ varia continuamente lungo la coesistenza.
\end{frame}

% 27
\begin{frame}{La pressione nella fase mista di Gibbs: primo passo}
\small
In ogni fase, a $T=0$:
\[
dP_i=n_B^{(i)}d\mu_B+n_Q^{(i)}d\mu_Q.
\]
Lungo la linea di coesistenza:
\[
P_{\mathrm I}=P_{\mathrm{II}}=P,
\]
quindi le due espressioni del differenziale sono uguali:
\[
\begin{aligned}
dP&=n_B^{\mathrm I}d\mu_B+n_Q^{\mathrm I}d\mu_Q,\\
dP&=n_B^{\mathrm{II}}d\mu_B+n_Q^{\mathrm{II}}d\mu_Q.
\end{aligned}
\]
Moltiplichiamo la prima equazione per $(1-\chi)$,
la seconda per $\chi$, poi sommiamo.
\end{frame}

% 28
\begin{frame}{La pressione nella fase mista di Gibbs: cancellazione}
\footnotesize
Sommando i differenziali con pesi $(1-\chi)$ e $\chi$:
\[
\begin{aligned}
dP={}&[(1-\chi)n_B^{\mathrm I}+\chi n_B^{\mathrm{II}}]d\mu_B\\
&+[(1-\chi)n_Q^{\mathrm I}+\chi n_Q^{\mathrm{II}}]d\mu_Q.
\end{aligned}
\]
Il secondo coefficiente e' nullo per neutralita' globale, mentre
il primo e' $n_B^{\rm tot}$. Dunque
\[
dP=n_B^{\rm tot}\,d\mu_B,
\qquad
\frac{dP}{d\mu_B}=n_B^{\rm tot}>0.
\]
\medskip
Il risultato vale \emph{lungo il percorso di coesistenza neutro}.
A differenza del plateau Maxwell, la pressione puo' crescere
al variare della composizione e della densita'.
\end{frame}

% 29
\begin{frame}{Perche' Gibbs permette una fase mista estesa?}
\small
\begin{columns}[T,onlytextwidth]
\begin{column}{0.49\textwidth}
\textbf{Maxwell}
\[
P=P_t\quad\text{nel tratto di coesistenza}.
\]
Nella soluzione TOV il passaggio avviene a un raggio $r_t$,
con salto di $\varepsilon$ e $P$ continuo.
\end{column}
\begin{column}{0.49\textwidth}
\textbf{Gibbs}
\[
\frac{dP}{d\mu_B}=n_B^{\rm tot}>0.
\]
Un gradiente di pressione e' possibile
nella regione mista; $\chi(r)$ puo' variare
attraverso uno strato del core.
\end{column}
\end{columns}
\medskip
La realizzazione effettiva dipende anche dalle energie di superficie,
Coulomb e screening, non incluse nella costruzione bulk.
\end{frame}

% 30
\begin{frame}{Maxwell e Gibbs: il ruolo della neutralita'}
\small
\begin{columns}[T,onlytextwidth]
\begin{column}{0.49\textwidth}
\textbf{Maxwell}
\begin{itemize}
\item Neutralita' \emph{locale} di entrambe le fasi.
\item Un potenziale indipendente dopo i vincoli.
\item Plateau a pressione $P_t$.
\item Salto di densita' a un'interfaccia idealmente netta.
\end{itemize}
\end{column}
\begin{column}{0.49\textwidth}
\textbf{Gibbs bulk}
\begin{itemize}
\item Neutralita' \emph{globale}; le singole fasi possono essere cariche.
\item Due potenziali: $\mu_B,\mu_Q$.
\item Pressione in generale crescente lungo la coesistenza.
\item Possibile strato misto con $\chi(r)$ variabile.
\end{itemize}
\end{column}
\end{columns}
\medskip
Non sono necessariamente due microfisiche diverse: cambiano
le ipotesi sulla distribuzione della carica e sul costo delle interfacce.
\end{frame}

% 31
\begin{frame}{Fase mista e geometrie ``pasta''}
\small
Se le due fasi sono localmente cariche ma globalmente neutre,
possono formare domini di una fase immersi nell'altra.
\medskip
Possibili geometrie:
\begin{itemize}
\item gocce quasi sferiche;
\item cilindri (``spaghetti'');
\item strati (``lasagne'');
\item strutture inverse, con ruoli delle fasi scambiati.
\end{itemize}
\medskip
La geometria e la scala dei domini dipendono dalla competizione fra
\[
E_S\quad\text{(superficie)},
\qquad E_C\quad\text{(Coulomb)}.
\]
Studiamo l'esempio del testo: gocce sferiche di raggio $R$.
\end{frame}

% 32
\begin{frame}{Energia superficiale: equazione (9.14)}
\small
Sia $\sigma$ la tensione superficiale fra le fasi
(energia per unita' di superficie).
\medskip
Per una goccia sferica di raggio $R$:
\[
S=4\pi R^2.
\]
L'energia associata all'interfaccia e'
\[
E_S=\sigma S=4\pi\sigma R^2.
\]
\begin{columns}[T,onlytextwidth]
\begin{column}{0.49\textwidth}
$\sigma$ misura il costo energetico della superficie di separazione.
\end{column}
\begin{column}{0.49\textwidth}
A parita' di $\sigma$:
\[
E_S\propto R^2.
\]
\end{column}
\end{columns}
\end{frame}

% 33
\begin{frame}{Contrasto di carica della goccia: equazione (9.15)}
\small
Indichiamo con
\[
\Delta n_Q=n_Q^{\mathrm I}-n_Q^{\mathrm{II}}
\]
la differenza delle densita' di carica (in unita' di $e$ per volume).
\medskip
Il volume della goccia e'
\[
V_d=\frac{4\pi}{3}R^3.
\]
Il contrasto di carica, in unita' della carica elementare, e'
\[
\Delta Z=V_d\Delta n_Q
=\frac{4\pi}{3}R^3\Delta n_Q.
\]
La carica elettrica associata al contrasto e' $e\Delta Z$;
la carica viene compensata dal mezzo circostante.
\end{frame}

% 34
\begin{frame}{Energia coulombiana: equazione (9.16)}
\small
Per una distribuzione sferica uniformemente carica,
nell'approssimazione elettrostatica del libro:
\[
E_C=\frac35\frac{(e\Delta Z)^2}{4\pi R}
=\frac{3\alpha_{\rm EM}}{5R}(\Delta Z)^2,
\]
con $\alpha_{\rm EM}=e^2/(4\pi)$ in unita' naturali.
\medskip
Usando $\Delta Z=(4\pi/3)R^3\Delta n_Q$:
\[
\begin{aligned}
E_C&=\frac{3\alpha_{\rm EM}}{5R}
\left(\frac{4\pi}{3}R^3\Delta n_Q\right)^2\\
&=\frac{16\pi^2\alpha_{\rm EM}}{15}
(\Delta n_Q)^2R^5.
\end{aligned}
\]
\end{frame}

% 35
\begin{frame}{Due energie e due diverse dipendenze da $R$}
\small
\[
E_S=4\pi\sigma R^2,
\qquad
E_C=\frac{16\pi^2\alpha_{\rm EM}}{15}
(\Delta n_Q)^2R^5.
\]
La prima cresce come $R^2$, la seconda come $R^5$.
\medskip
\textbf{Attenzione:} minimizzare $E_S+E_C$ per \emph{una sola goccia}
non determina il raggio ottimale: entrambi i termini crescono con $R$.
\medskip
Occorre invece fissare la frazione volumetrica della fase dispersa
e minimizzare l'energia \emph{per unita' di volume totale}.
\end{frame}

% 36
\begin{frame}{Quante gocce servono a frazione volumetrica fissata?}
\small
Sia $\chi$ la frazione di volume occupata dalle gocce.
\medskip
Ogni goccia ha volume
\[
V_d=\frac{4\pi}{3}R^3.
\]
Il numero di gocce per unita' di volume totale e'
\[
\mathcal N=\frac{\chi}{V_d}
=\frac{3\chi}{4\pi R^3}.
\]
Se $R$ aumenta mantenendo $\chi$ costante, diminuisce il numero di gocce:
\[
\mathcal N\propto R^{-3}.
\]
Questo fattore ribalta la dipendenza dell'energia superficiale
\emph{per unita' di volume}.
\end{frame}

% 37
\begin{frame}{Energia superficiale per unita' di volume}
\small
Moltiplichiamo l'energia di una goccia per la loro densita' numerica:
\[
\mathcal E_S=\mathcal N E_S.
\]
Sostituendo:
\[
\begin{aligned}
\mathcal E_S
&=\frac{3\chi}{4\pi R^3}\,(4\pi\sigma R^2)\\
&=\frac{3\chi\sigma}{R}.
\end{aligned}
\]
Quindi
\[
\mathcal E_S\propto R^{-1}.
\]
A volume totale fissato, poche gocce grandi
possiedono meno superficie complessiva di molte gocce piccole.
\end{frame}

% 38
\begin{frame}{Energia coulombiana per unita' di volume}
\small
Analogamente,
\[
\mathcal E_C=\mathcal N E_C.
\]
Sostituiamo le espressioni precedenti:
\[
\begin{aligned}
\mathcal E_C
&=\frac{3\chi}{4\pi R^3}
\frac{16\pi^2\alpha_{\rm EM}}{15}
(\Delta n_Q)^2R^5\\
&=\frac{4\pi}{5}\chi\alpha_{\rm EM}
(\Delta n_Q)^2R^2.
\end{aligned}
\]
Quindi
\[
\mathcal E_C\propto R^2.
\]
L'energia superficiale favorisce gocce grandi,
quella coulombiana gocce piccole.
\end{frame}

% 39
\begin{frame}{Raggio ottimale delle gocce: derivazione}
\footnotesize
L'energia addizionale \emph{per volume totale} e'
\[
\mathcal E(R)=\frac{3\chi\sigma}{R}
+\frac{4\pi}{5}\chi\alpha_{\rm EM}(\Delta n_Q)^2R^2.
\]
La condizione di minimo e'
\[
\begin{aligned}
0=\frac{d\mathcal E}{dR}
&=-\frac{3\chi\sigma}{R^2}
+\frac{8\pi}{5}\chi\alpha_{\rm EM}(\Delta n_Q)^2R.
\end{aligned}
\]
Riordinando e cancellando $\chi$:
\[
R^3=\frac{15\sigma}
{8\pi\alpha_{\rm EM}(\Delta n_Q)^2}.
\]
Da cui l'equazione (9.17):
\[
R=\left[\frac{15\sigma}
{8\pi\alpha_{\rm EM}(\Delta n_Q)^2}\right]^{1/3}.
\]
\end{frame}

% 40
\begin{frame}{Verifica del minimo: $E_S=2E_C$}
\small
Dalla condizione di stazionarieta'
\[
\frac{3\chi\sigma}{R^2}
=\frac{8\pi}{5}\chi\alpha_{\rm EM}(\Delta n_Q)^2R,
\]
moltiplichiamo entrambi i membri per $R$:
\[
\frac{3\chi\sigma}{R}
=2\left[\frac{4\pi}{5}\chi\alpha_{\rm EM}
(\Delta n_Q)^2R^2\right].
\]
Riconosciamo quindi
\[
\mathcal E_S=2\mathcal E_C,
\qquad E_S=2E_C.
\]
Questa relazione dipende dalle specifiche leggi di scala
adottate per gocce sferiche nel modello.
\end{frame}

% 41
\begin{frame}{Stima del raggio: equazione (9.17)}
\small
Ripristinando le unita' MeV e fm, il libro riporta
\[
\begin{aligned}
R\simeq{}&5.7\,\mathrm{fm}
\left(\frac{\sigma}{10\,\mathrm{MeV\,fm^{-2}}}\right)^{1/3}\\
&\times\left(\frac{|\Delta n_Q|}{0.15\,\mathrm{fm^{-3}}}\right)^{-2/3}.
\end{aligned}
\]
\begin{columns}[T,onlytextwidth]
\begin{column}{0.49\textwidth}
\[
R\propto\sigma^{1/3}.
\]
Una superficie piu' costosa favorisce gocce piu' grandi.
\end{column}
\begin{column}{0.49\textwidth}
\[
R\propto|\Delta n_Q|^{-2/3}.
\]
Un contrasto di carica maggiore favorisce gocce piu' piccole.
\end{column}
\end{columns}
\end{frame}

% 42
\begin{frame}{Quando la fase mista e' energeticamente favorita?}
\small
La costruzione di Gibbs bulk trascura il costo dei domini.
Con superficie e Coulomb occorre confrontare
\[
\varepsilon_{\rm mix}
=\varepsilon_{\rm bulk}+\mathcal E_S+\mathcal E_C+\cdots
\]
con l'energia della separazione macroscopica in fasi quasi neutre.
\begin{columns}[T,onlytextwidth]
\begin{column}{0.49\textwidth}
\textbf{$\sigma$ piccola}
\medskip
La separazione in domini caricati puo' essere favorevole:
fase mista e ``pasta'' nel core.
\end{column}
\begin{column}{0.49\textwidth}
\textbf{$\sigma$ grande}
\medskip
Molte interfacce costano troppo:
configurazioni vicine al limite di Maxwell.
\end{column}
\end{columns}
\end{frame}

% 43
\begin{frame}{Limiti del modello di goccia uniforme}
\small
\begin{itemize}
\item La tensione superficiale adroni--quark $\sigma$ non e' nota con precisione.
\item Le cariche si possono redistribuire per schermatura elettrica.
\item Forme cilindriche, planari o inverse possono sostituire le gocce.
\item Le lunghezze di screening vanno confrontate con la scala $R$.
\end{itemize}
\medskip
Il libro cita stime indicative:
\[
\sigma\sim10\,\mathrm{MeV\,fm^{-2}},
\qquad
\sigma_{\rm crit}\sim70\,\mathrm{MeV\,fm^{-2}}.
\]
\footnotesize
Sono valori dipendenti dal modello, non costanti universali della QCD.
Per grandi tensioni il guadagno bulk da neutralita' globale
puo' non compensare il costo di superficie e Coulomb.
\end{frame}

% 44
\begin{frame}{Come cambia la struttura del core stellare?}
\small
\begin{columns}[T,onlytextwidth]
\begin{column}{0.49\textwidth}
\textbf{Maxwell: interfaccia netta}
\[
P(r_t^-)=P(r_t^+)=P_t,
\]
\[
\varepsilon(r_t^-)\ne\varepsilon(r_t^+).
\]
Un core di fase II puo' essere separato
bruscamente dall'involucro adronico.
\end{column}
\begin{column}{0.49\textwidth}
\textbf{Gibbs: fase mista estesa}
\[
0<\chi(r)<1.
\]
Le due fasi possono coesistere
in una regione radiale, con pressione
variabile e possibili strutture pasta.
\end{column}
\end{columns}
\medskip
Una transizione sufficientemente intensa puo' ammorbidire l'EOS,
modificare $M(R)$ e, in certi modelli, produrre rami di stabilita' distinti.
\end{frame}

% 45
\begin{frame}{Riepilogo: equilibrio e geometria della transizione}
\footnotesize
\textbf{Maxwell (una variabile chimica indipendente):}
\[
P_{\mathrm I}(\mu_t)=P_{\mathrm{II}}(\mu_t),
\qquad \Delta\varepsilon=\mu_t\Delta n_B.
\]
\textbf{Gibbs (due cariche conservate):}
\[
P_{\mathrm I}(\mu_B,\mu_Q)
=P_{\mathrm{II}}(\mu_B,\mu_Q).
\]
\textbf{Neutralita' della miscela:}
\[
(1-\chi)n_Q^{\mathrm I}+\chi n_Q^{\mathrm{II}}=0.
\]
\textbf{Scala delle gocce sferiche:}
\[
R=\left[\frac{15\sigma}
{8\pi\alpha_{\rm EM}(\Delta n_Q)^2}\right]^{1/3}.
\]
La topologia del core dipende dall'EOS, dai vincoli di carica
e dal costo delle interfacce.
\end{frame}

\section{Terza famiglia di stelle compatte}

\subsection{9.3.1 Twin Stars}

\begin{frame}{Terza famiglia: problema fisico}
\small
Una transizione del primo ordine puo' modificare la topologia
 delle soluzioni stabili delle equazioni TOV.
\[
\text{stabile}\;\longrightarrow\;\text{instabile}
\;\longrightarrow\;\text{nuovamente stabile}.
\]
\begin{itemize}
\item Prima famiglia: nane bianche.
\item Seconda famiglia: stelle di neutroni.
\item Terza famiglia: eventuale nuovo ramo stabile ad alta densita'.
\end{itemize}
Una terza famiglia non richiede necessariamente quark liberi:
richiede una EOS che consenta una nuova sequenza stabile.
\end{frame}

\begin{frame}{Sequenze di equilibrio e modi radiali}
\small
\begin{columns}[T,onlytextwidth]
\begin{column}{0.49\textwidth}
Le configurazioni TOV formano una sequenza al variare
 della pressione centrale $P_c$:
\[
M=M(P_c),\qquad R=R(P_c).
\]
La stabilita' radiale si analizza mediante i modi normali:
\[
\omega_0^2,\omega_1^2,\omega_2^2,\ldots
\]
\end{column}
\begin{column}{0.49\textwidth}
\begin{itemize}
\item $\omega_n^2>0$: oscillazione stabile del modo $n$.
\item $\omega_n^2<0$: modo instabile.
\item $\omega_n^2=0$: soglia di stabilita'.
\end{itemize}
Infatti, se $\omega_n^2=-\gamma_n^2$,
\[
\xi_n(t)\propto e^{\gamma_nt}.
\]
Per stabilita' radiale occorre che \emph{tutti} i modi siano stabili.
\end{column}
\end{columns}
\end{frame}

\begin{frame}{Figura 9.6: stabilita' e terza famiglia}
\small
\begin{figure}[!ht]\includegraphics[trim={0cm 0cm 0 0},clip,keepaspectratio,width=0.95\textwidth]{Schaffner/stabilityregions_thirdfamily}\label{fig:Schaffner_stabilityregions_thirdfamily}\end{figure}
\footnotesize Le caselle chiare rappresentano modi stabili; quelle scure modi instabili.
\end{frame}

\begin{frame}{Come leggere le sequenze della Figura 9.6}
\small
\begin{columns}[T,onlytextwidth]
\begin{column}{0.49\textwidth}
\begin{itemize}
\item Nane bianche $\to A$: stabili.
\item $A\to B$: instabili dopo il limite di Chandrasekhar.
\item $B\to C$: ramo stabile delle stelle di neutroni.
\end{itemize}
\end{column}
\begin{column}{0.49\textwidth}
\begin{itemize}
\item $C\to D$: intervallo instabile.
\item $D\to E$: nuova regione stabile: terza famiglia.
\item Oltre $E$: nuova perdita di stabilita'.
\end{itemize}
\medskip
Gli altri estremi possono modificare il numero di modi instabili
senza necessariamente ristabilire il modo fondamentale.
\end{column}
\end{columns}
\end{frame}

\begin{frame}{Estremi della curva e cambiamenti di stabilita'}
\small
Per una sequenza TOV a un parametro, gli estremi della massa
possono segnalare l'attraversamento di un modo a frequenza nulla.
\[
\frac{dM}{dP_c}=0.
\]
\begin{itemize}
\item Bisogna seguire il verso di percorrenza al crescere di $P_c$.
\item Un massimo o un minimo di $M(R)$ non identifica \emph{da solo}
quale modo cambia stabilita'.
\item Un modo superiore che si ristabilizza non compensa
un modo fondamentale ancora instabile.
\end{itemize}
I salti di densita' del primo ordine richiedono anche di controllare
la stabilita' alla comparsa della nuova fase.
\end{frame}

\begin{frame}{Twin stars: stessa massa, raggi differenti}
\small
\begin{columns}[T,onlytextwidth]
\begin{column}{0.49\textwidth}
Due configurazioni \emph{stabili} appartenenti a rami differenti
possono soddisfare
\[
M_1\simeq M_2,\qquad R_1
eq R_2.
\]
Esempio puramente illustrativo:
\[
M_1=M_2=1.9M_\odot,
\]
\[
R_1=13\,\mathrm{km},\quad R_2=10\,\mathrm{km}.
\]
\end{column}
\begin{column}{0.49\textwidth}
\begin{itemize}
\item Stella meno compatta: materia prevalentemente adronica
 o piccolo core della nuova fase.
\item Stella piu' compatta: core denso piu' esteso.
\end{itemize}
\medskip
Una misura di raggi diversi a masse simili sarebbe interessante,
ma non dimostrerebbe, da sola, il deconfinamento.
\end{column}
\end{columns}
\end{frame}

\begin{frame}{Terza famiglia, stabilita' e contenuto microscopico}
\small
Le TOV dipendono dall'equazione di stato
\[
P=P(\varepsilon),
\]
non dalla particolare interpretazione microscopica adottata.
\medskip
Una nuova fase puo' essere materia di quark o un'altra
fase capace di produrre un forte cambiamento della EOS.
\medskip
\textbf{Due questioni distinte:}
\begin{itemize}
\item esiste un \emph{secondo ramo stabile} della sequenza $M(R)$?
\item a pari massa gravitazionale, le due configurazioni
hanno lo stesso numero barionico? Non necessariamente.
\end{itemize}
\end{frame}

\subsection{9.3.2 The Seidov Criterion}

\begin{frame}{Il problema di Seidov}
\small
Transizione di Maxwell alla pressione $P_0$:
\[
\varepsilon_1\longrightarrow\varepsilon_2,
\qquad \varepsilon_2>\varepsilon_1,
\qquad \Delta\varepsilon=\varepsilon_2-\varepsilon_1.
\]
Appena $P_c$ supera $P_0$ appare un core della fase II di raggio $r_n$.
\medskip
\textbf{Domanda:} quando la comparsa di un core infinitesimo
rende instabile una sequenza precedentemente stabile?
\medskip
Il criterio cercato e' locale alla \emph{soglia di formazione}
del core, non alla successiva evoluzione del ramo ibrido.
\end{frame}

\begin{frame}{TOV con $G$ esplicita e $c=1$}
\small
\begin{columns}[T,onlytextwidth]
\begin{column}{0.49\textwidth}
Equilibrio idrostatico:
\[
\frac{dP}{dr}=-G\frac{(P+\varepsilon)
[m(r)+4\pi r^3P]}{r[r-2Gm(r)]}.
\]
Massa interna:
\[
\frac{dm}{dr}=4\pi r^2\varepsilon.
\]
\end{column}
\begin{column}{0.49\textwidth}
\begin{itemize}
\item $r$: coordinata radiale areale.
\item $m(r)$: massa gravitazionale interna.
\item $P(r)$, $\varepsilon(r)$: pressione e densita' di energia.
\item $P_0$: pressione della transizione.
\item $\varepsilon_1,\varepsilon_2$: densita' ai due lati della transizione.
\end{itemize}
\medskip
Studiamo il limite $r\to0$.
\end{column}
\end{columns}
\end{frame}

\begin{frame}{Configurazione di riferimento: massa centrale}
\small
Prima della transizione:
\[
P_c=P_0,\qquad\varepsilon_c=\varepsilon_1.
\]
Per $r$ piccolo prendiamo $\varepsilon(r)\simeq\varepsilon_1$.
\[
\frac{dm_-}{dr}\simeq4\pi r^2\varepsilon_1.
\]
Integrando con $m_-(0)=0$:
\[
\begin{aligned}
m_-(r)&=4\pi\varepsilon_1\int_0^r r'^2\,dr'\\
&=\frac{4\pi}{3}\varepsilon_1r^3.
\end{aligned}
\]
\end{frame}

\begin{frame}{Ordine del fattore di Schwarzschild}
\small
Dallo sviluppo della massa:
\[
\frac{2Gm_-(r)}r\simeq
\frac{8\pi G}{3}\varepsilon_1r^2=O(r^2).
\]
Quindi
\[
\frac{1}{r[r-2Gm_-(r)]}
=\frac{1}{r^2}\left[1+O(r^2)\right].
\]
Al primo ordine non nullo nel gradiente $dP/dr$ possiamo
sostituire il denominatore con $r^2$.
\medskip
\textbf{Attenzione:} i termini relativistici che coinvolgono $P$
nel numeratore delle TOV non vengono eliminati.
\end{frame}

\begin{frame}{Gradiente di pressione prima della transizione}
\small
Sostituiamo $m_-(r)\simeq(4\pi/3)\varepsilon_1r^3$:
\[
\begin{aligned}
\frac{dP_-}{dr}&\simeq-G\frac{P_0+\varepsilon_1}{r^2}
\left(4\pi r^3P_0+\frac{4\pi}{3}r^3\varepsilon_1\right)\\
&=-\frac{4\pi G}{3}(P_0+\varepsilon_1)
(3P_0+\varepsilon_1)r.
\end{aligned}
\]
Il gradiente si annulla linearmente per $r\to0$, come richiesto
dalla regolarita' al centro.
\end{frame}

\begin{frame}{Profilo di pressione della soluzione $-$}
\small
Integrando lo sviluppo del gradiente:
\[
\begin{aligned}
P_-(r)-P_0
&=-\frac{4\pi G}{3}(P_0+\varepsilon_1)
(3P_0+\varepsilon_1)\int_0^r r'\,dr'\\
&=-\frac{2\pi G}{3}(P_0+\varepsilon_1)
(3P_0+\varepsilon_1)r^2.
\end{aligned}
\]
Dunque
\[
P_-(r)\simeq P_0-\frac{2\pi G}{3}
(P_0+\varepsilon_1)(3P_0+\varepsilon_1)r^2.
\]
\end{frame}

\begin{frame}{Precisazione sulle formule del testo}
\small
La sostituzione diretta nelle TOV fornisce
\[
(P_0+\varepsilon_1)(3P_0+\varepsilon_1),
\]
non il fattore
\[
(P_0+\varepsilon_1)(P_0+3\varepsilon_1)
\]
che compare in alcuni passaggi intermedi del testo allegato.
\medskip
Useremo gli sviluppi \emph{algebricamente coerenti} delle TOV:
essi riproducono il coefficiente finale di Seidov.
\medskip
Le equazioni originali sono riportate nella fonte come
(9.18)--(9.35); i numeri non sono riprodotti sulle slide.
\end{frame}

\begin{frame}{Configurazione $+$: nasce un piccolo core}
\small
Superiamo appena la pressione di transizione:
\[
P_c=P_0+\delta,\qquad\delta>0.
\]
Compare un core della fase ad alta densita' con raggio $r_n\ll R$:
\[
\varepsilon(r<r_n)\simeq\varepsilon_2.
\]
La pressione al bordo del core e' continua:
\[
P_+(r_n)=P_0.
\]
All'esterno del core si trova ancora la fase di bassa densita'
$\varepsilon_1$.
\end{frame}

\begin{frame}{Massa del piccolo core}
\small
Nel core $\varepsilon\simeq\varepsilon_2$:
\[
\begin{aligned}
m_+(r_n)
&=4\pi\int_0^{r_n}\varepsilon_2r^2\,dr\\
&=\frac{4\pi}{3}\varepsilon_2r_n^3.
\end{aligned}
\]
Quindi
\[
{m_+(r_n)\propto r_n^3.}
\]
La massa aggiuntiva rispetto alla fase I nello stesso volume e'
\[
\Delta m\simeq\frac{4\pi}{3}
(\varepsilon_2-\varepsilon_1)r_n^3.
\]
\end{frame}

\begin{frame}{Pressione all'interno del nuovo core}
\small
Sostituendo $P\simeq P_0$ e $\varepsilon\simeq\varepsilon_2$
nello sviluppo centrale delle TOV:
\[
\frac{dP_+}{dr}\simeq-\frac{4\pi G}{3}
(P_0+\varepsilon_2)(3P_0+\varepsilon_2)r.
\]
Integrando con $P_+(0)=P_0+\delta$:
\[
{
P_+(r)\simeq P_0+\delta-
\frac{2\pi G}{3}(P_0+\varepsilon_2)
(3P_0+\varepsilon_2)r^2.}
\]
\end{frame}

\begin{frame}{Raggio del core e incremento $\delta$}
\small
Al bordo $P_+(r_n)=P_0$, quindi:
\[
\begin{aligned}
P_0&=P_0+\delta-
\frac{2\pi G}{3}(P_0+\varepsilon_2)
(3P_0+\varepsilon_2)r_n^2,\\
\delta&=\frac{2\pi G}{3}(P_0+\varepsilon_2)
(3P_0+\varepsilon_2)r_n^2.
\end{aligned}
\]
Ne consegue
\[
{r_n\propto\sqrt{\delta}.}
\]
Il core nasce con raggio nullo e cresce quando $P_c$ supera $P_0$.
\end{frame}

\begin{frame}{Gradiente appena all'esterno del core}
\small
Al bordo esterno $r_n^+$ il fluido e' nella fase I:
\[
\varepsilon_{\rm locale}=\varepsilon_1.
\]
La massa interna e' pero' quella del core denso:
\[
m(r_n)=\frac{4\pi}{3}\varepsilon_2r_n^3.
\]
Percio' le TOV danno
\[
{
\left.\frac{dP_+}{dr}\right|_{r_n^+}
\simeq-\frac{4\pi G}{3}(P_0+\varepsilon_1)
(3P_0+\varepsilon_2)r_n.}
\]
La distinzione fra densita' locale e massa interna e' essenziale.
\end{frame}

\begin{frame}{Confronto dei due profili di pressione}
\small
Definiamo la perturbazione nel mantello di bassa densita':
\[
\Pi(r)=P_+(r)-P_-(r).
\]
Per il riferimento:
\[
\left.\frac{dP_-}{dr}\right|_{r_n}
=-\frac{4\pi G}{3}(P_0+\varepsilon_1)
(3P_0+\varepsilon_1)r_n.
\]
Per il nuovo core:
\[
\left.\frac{dP_+}{dr}\right|_{r_n^+}
=-\frac{4\pi G}{3}(P_0+\varepsilon_1)
(3P_0+\varepsilon_2)r_n.
\]
\end{frame}

\begin{frame}{Linearizzazione: $\Pi(r)=A+B/r$}
\small
La soluzione perturbativa asintotica nel mantello esterno assume la forma
\[
{\Pi(r)=A+\frac{B}{r}.}
\]
\begin{itemize}
\item $A$: spostamento del profilo esterno regolare al limite $r_n\to0$.
\item $B/r$: effetto della massa in eccesso del core.
\item $A$ e $B$ si determinano dalla pressione e dal gradiente a $r_n$.
\end{itemize}
La forma e' valida all'ordine dominante nello sviluppo per un core
infinitesimo: non rappresenta una soluzione TOV globale esatta.
\end{frame}

\begin{frame}{Prima condizione al bordo: la pressione}
\small
Dal profilo della stella di riferimento:
\[
P_-(r_n)=P_0-
\frac{2\pi G}{3}(P_0+\varepsilon_1)
(3P_0+\varepsilon_1)r_n^2.
\]
Poiche' $P_+(r_n)=P_0$:
\[
\Pi(r_n)=\frac{2\pi G}{3}(P_0+\varepsilon_1)
(3P_0+\varepsilon_1)r_n^2.
\]
Dalla forma asintotica:
\[
{A+\frac{B}{r_n}=
\frac{2\pi G}{3}(P_0+\varepsilon_1)
(3P_0+\varepsilon_1)r_n^2.}
\]
\end{frame}

\begin{frame}{Seconda condizione: il gradiente}
\small
Sottraiamo i due gradienti al bordo:
\[
\begin{aligned}
\left.\frac{d\Pi}{dr}\right|_{r_n^+}
&=\left.\frac{dP_+}{dr}\right|_{r_n^+}
-\left.\frac{dP_-}{dr}\right|_{r_n}\\
&=-\frac{4\pi G}{3}(P_0+\varepsilon_1)
(\varepsilon_2-\varepsilon_1)r_n.
\end{aligned}
\]
Dalla soluzione $\Pi(r)=A+B/r$:
\[
{-\frac{B}{r_n^2}=
-\frac{4\pi G}{3}(P_0+\varepsilon_1)
(\varepsilon_2-\varepsilon_1)r_n.}
\]
\end{frame}

\begin{frame}{Determinazione del coefficiente $B$}
\small
Partiamo dalla seconda condizione al bordo:
\[
-\frac{B}{r_n^2}=-\frac{4\pi G}{3}
(P_0+\varepsilon_1)(\varepsilon_2-\varepsilon_1)r_n.
\]
Moltiplichiamo per $-r_n^2$:
\[
{B=\frac{4\pi G}{3}
(P_0+\varepsilon_1)(\varepsilon_2-\varepsilon_1)r_n^3.}
\]
Pertanto il termine singolare $B/r$ e' di ordine $r_n^3/r$,
come ci si aspetta per l'effetto della piccola massa aggiuntiva.
\end{frame}

\begin{frame}{Determinazione di $A$: sostituzione}
\small
Dalla prima condizione:
\[
A=\frac{2\pi G}{3}(P_0+\varepsilon_1)
(3P_0+\varepsilon_1)r_n^2-\frac{B}{r_n}.
\]
Sostituiamo $B$:
\[
\begin{aligned}
A={}&\frac{2\pi G}{3}(P_0+\varepsilon_1)r_n^2
\big[(3P_0+\varepsilon_1)\\
&\hspace{31mm}-2(\varepsilon_2-\varepsilon_1)\big].
\end{aligned}
\]
\end{frame}

\begin{frame}{Determinazione di $A$: semplificazione}
\small
Semplifichiamo la parentesi:
\[
\begin{aligned}
3P_0+\varepsilon_1-2(\varepsilon_2-\varepsilon_1)
&=3P_0+\varepsilon_1-2\varepsilon_2+2\varepsilon_1\\
&=3P_0+3\varepsilon_1-2\varepsilon_2.
\end{aligned}
\]
Quindi
\[
{
A=\frac{2\pi G}{3}(P_0+\varepsilon_1)
(3P_0+3\varepsilon_1-2\varepsilon_2)r_n^2.}
\]
Questa e' l'espressione del testo associata all'equazione (9.31).
\end{frame}

\begin{frame}{Eliminiamo $r_n$ mediante $\delta$}
\small
Ricordiamo
\[
\delta=\frac{2\pi G}{3}
(P_0+\varepsilon_2)(3P_0+\varepsilon_2)r_n^2.
\]
Da cui
\[
r_n^2=\frac{3\delta}{2\pi G
(P_0+\varepsilon_2)(3P_0+\varepsilon_2)}.
\]
Sostituendo in $A$:
\[
{\frac{A}{\delta}=
\frac{(P_0+\varepsilon_1)
(3P_0+3\varepsilon_1-2\varepsilon_2)}
{(P_0+\varepsilon_2)(3P_0+\varepsilon_2)}.}
\]
\end{frame}

\begin{frame}{Derivata della massa e cambio di segno}
\small
Nel limite di comparsa del core, il testo collega
la risposta del ramo con e senza core mediante
\[
{
\frac{dM_+}{dP_c}=\frac{A}{\delta}
\frac{dM_-}{dP_c}.}
\]
Se il ramo iniziale e' stabile e $dM_-/dP_c>0$,
il segno della nuova derivata dipende da $A/\delta$.
\medskip
Per $P_0>0$, $\varepsilon_{1,2}>0$, il denominatore
 di $A/\delta$ e' positivo: resta da studiare
\[
3P_0+3\varepsilon_1-2\varepsilon_2.
\]
\end{frame}

\begin{frame}{Criterio di Seidov: forma fondamentale}
\small
L'instabilita' immediata richiede
\[
3P_0+3\varepsilon_1-2\varepsilon_2<0.
\]
Quindi
\[
2\varepsilon_2>3(P_0+\varepsilon_1).
\]
Dividendo per $2\varepsilon_1$:
\[
{
\frac{\varepsilon_2}{\varepsilon_1}>
\frac32\left(1+\frac{P_0}{\varepsilon_1}\right).}
\]
\textbf{Interpretazione:} un salto energetico sufficientemente grande
destabilizza la stella quando compare un core infinitesimo.
\end{frame}

\begin{frame}{Criterio di Seidov: salto relativo}
\small
Definiamo
\[
\Delta\varepsilon=\varepsilon_2-\varepsilon_1.
\]
Poiche' $\varepsilon_2=\varepsilon_1+\Delta\varepsilon$,
\[
1+\frac{\Delta\varepsilon}{\varepsilon_1}>
\frac32\left(1+\frac{P_0}{\varepsilon_1}\right).
\]
Sottraendo 1:
\[
{
\frac{\Delta\varepsilon}{\varepsilon_1}>
\frac12+\frac32\frac{P_0}{\varepsilon_1}.}
\]
Questa e' la forma impiegata nella classificazione delle stelle ibride.
\end{frame}

\begin{frame}{Limite newtoniano del criterio}
\small
Se la pressione alla transizione e' piccola
rispetto alla densita' di energia,
\[
\frac{P_0}{\varepsilon_1}\ll1,
\]
allora
\[
{\frac{\varepsilon_2}{\varepsilon_1}>\frac32.}
\]
Oppure
\[
{\frac{\Delta\varepsilon}{\varepsilon_1}>\frac12.}
\]
Un salto superiore al $50\%$ e' quindi la soglia
nell'approssimazione newtoniana.
\end{frame}

\begin{frame}{Esempio numerico per il salto critico}
\small
\begin{columns}[T,onlytextwidth]
\begin{column}{0.49\textwidth}
Prendiamo
\[
\varepsilon_1=500\,\mathrm{MeV\,fm^{-3}},
\]
\[
P_0=100\,\mathrm{MeV\,fm^{-3}}.
\]
Allora
\[
\frac{P_0}{\varepsilon_1}=0.2.
\]
\end{column}
\begin{column}{0.49\textwidth}
Il criterio richiede
\[
\frac{\varepsilon_2}{\varepsilon_1}>
\frac32(1+0.2)=1.8.
\]
Quindi
\[
\varepsilon_2>900\,\mathrm{MeV\,fm^{-3}},
\]
\[
\Delta\varepsilon>400\,\mathrm{MeV\,fm^{-3}}.
\]
Non dice se apparira' poi un nuovo ramo stabile.
\end{column}
\end{columns}
\end{frame}

\begin{frame}{Che cosa stabilisce il criterio e che cosa no}
\small
\begin{columns}[T,onlytextwidth]
\begin{column}{0.49\textwidth}
\textbf{Stabilisce}
\begin{itemize}
\item La stabilita' \emph{immediata} alla nascita della nuova fase.
\item Una soglia per il salto di energia a pressione fissata.
\item Il comportamento locale al limite $r_n\to0$.
\end{itemize}
\end{column}
\begin{column}{0.49\textwidth}
\textbf{Non stabilisce}
\begin{itemize}
\item La presenza di un secondo ramo stabile a pressioni elevate.
\item La rigidita' della fase densa.
\item L'esistenza osservabile di stelle gemelle.
\end{itemize}
\medskip
Per questi aspetti occorre studiare \emph{tutta} la sequenza TOV.
\end{column}
\end{columns}
\end{frame}

\subsection{9.3.3 Classification of Hybrid Stars}

\begin{frame}{Due parametri adimensionali della transizione}
\small
Alla pressione $P_{\rm trans}$ la densita' passa da
$\varepsilon_{\rm trans}$ a
$\varepsilon_{\rm trans}+\Delta\varepsilon$.
\[
x=\frac{P_{\rm trans}}{\varepsilon_{\rm trans}},
\qquad
y=\frac{\Delta\varepsilon}{\varepsilon_{\rm trans}}.
\]
Il criterio di Seidov diventa
\[
{y>\frac12+\frac32 x.}
\]
Nel piano $(x,y)$ la soglia di instabilita' immediata
 e' quindi una \emph{retta}.
\end{frame}

\begin{frame}{La parametrizzazione CSS della fase densa}
\small
Descriviamo la fase ad alta densita' tramite
\[
{P=s(\varepsilon-\varepsilon_*).}
\]
\begin{itemize}
\item $s=dP/d\varepsilon=c_s^2$ in unita' $c=1$.
\item $\varepsilon_*$: costante che fissa il punto a pressione nulla
  dell'estrapolazione lineare.
\item Causalita': $0<s\leq1$.
\end{itemize}
Il valore $s=1$ e' il limite di massima rigidita' causale
considerato nel testo.
\end{frame}

\begin{frame}{Determinazione della costante $\varepsilon_*$}
\small
All'inizio della fase ad alta densita':
\[
P_{\rm trans}=s(\varepsilon_{\rm trans}
+\Delta\varepsilon-\varepsilon_*).
\]
Risolvendo:
\[
{
\varepsilon_*=
\varepsilon_{\rm trans}+\Delta\varepsilon-
\frac{P_{\rm trans}}s.}
\]
Quindi
\[
P=P_{\rm trans}+s\big[\varepsilon-
(\varepsilon_{\rm trans}+\Delta\varepsilon)\big].
\]
La pressione e' continua, la densita' di energia salta.
\end{frame}

\begin{frame}{EOS completa per le equazioni TOV}
\footnotesize
\[
P(\varepsilon)=
\begin{cases}
P_{\rm low}(\varepsilon),
& \varepsilon<\varepsilon_{\rm trans},\\[1mm]
P_{\rm trans}
+s\,[\varepsilon-(\varepsilon_{\rm trans}+\Delta\varepsilon)],
& \varepsilon>\varepsilon_{\rm trans}+\Delta\varepsilon.
\end{cases}
\]
Nel tratto intermedio $\varepsilon_{\rm trans}<\varepsilon<
\varepsilon_{\rm trans}+\Delta\varepsilon$ la Maxwell
construction ha pressione costante $P_{\rm trans}$,
ma non genera uno strato stellare radiale esteso.
\medskip
Forma inversa per il core:
\[
{\varepsilon_Q(P)=\varepsilon_{\rm trans}
+\Delta\varepsilon+\frac{P-P_{\rm trans}}s.}
\]
\end{frame}

\begin{frame}{I due meccanismi antagonisti}
\small
\begin{columns}[T,onlytextwidth]
\begin{column}{0.49\textwidth}
\textbf{Salto energetico}
\[
\Delta\varepsilon\;\text{grande}.
\]
\begin{itemize}
\item Favorisce l'instabilita' iniziale.
\item Puo' interrompere la sequenza ordinaria.
\item La soglia locale e' data da Seidov.
\end{itemize}
\end{column}
\begin{column}{0.49\textwidth}
\textbf{Rigidita' della fase densa}
\[
s=\frac{dP}{d\varepsilon}\;\text{grande}.
\]
\begin{itemize}
\item Aumenta il sostegno pressorio.
\item Puo' favorire una nuova regione stabile.
\item Non garantisce da sola una terza famiglia.
\end{itemize}
\end{column}
\end{columns}
\end{frame}

\begin{frame}{Figura 9.7: classificazione A--D}
\small
\begin{figure}[!ht]\includegraphics[trim={0cm 0cm 0 0},clip,keepaspectratio,width=0.95\textwidth]{Schaffner/hybridclassification}\label{fig:Schaffner_hybridclassification}\end{figure}
\centering
\includegraphics[width=0.96\linewidth,height=0.75\textheight,keepaspectratio]{qcd93_fig97.png}

\footnotesize Asse verticale: salto relativo $\Delta\varepsilon/\varepsilon_{\rm trans}$;
orizzontale: $P_{\rm trans}/\varepsilon_{\rm trans}$.
\end{frame}

\begin{frame}{Regioni A e D: sopra la retta di Seidov}
\small
\begin{columns}[T,onlytextwidth]
\begin{column}{0.49\textwidth}
\textbf{Classe A}
\begin{itemize}
\item La transizione destabilizza immediatamente.
\item La pressione di transizione e' relativamente elevata.
\item Non emerge un ramo ibrido stabile.
\end{itemize}
\end{column}
\begin{column}{0.49\textwidth}
\textbf{Classe D}
\begin{itemize}
\item La transizione destabilizza immediatamente.
\item La pressione di transizione e' piu' bassa.
\item A pressioni centrali maggiori appare un
nuovo ramo stabile \emph{disconnesso}.
\end{itemize}
\end{column}
\end{columns}
\end{frame}

\begin{frame}{Regioni C e B: sotto la retta di Seidov}
\small
\begin{columns}[T,onlytextwidth]
\begin{column}{0.49\textwidth}
\textbf{Classe C}
\begin{itemize}
\item La comparsa del core non destabilizza subito.
\item Il ramo ibrido resta connesso a quello ordinario.
\item Non e' presente un secondo ramo stabile separato.
\end{itemize}
\end{column}
\begin{column}{0.49\textwidth}
\textbf{Classe B}
\begin{itemize}
\item Il core nasce senza instabilita' immediata.
\item La sequenza diventa instabile piu' avanti.
\item Riappare poi un ramo ibrido disconnesso.
\end{itemize}
\end{column}
\end{columns}
\end{frame}

\begin{frame}{Perche' la classe B e' importante?}
\small
La classe B si trova \emph{sotto} la retta di Seidov:
\[
\frac{\Delta\varepsilon}{\varepsilon_{\rm trans}}
<\frac12+\frac32\frac{P_{\rm trans}}{\varepsilon_{\rm trans}}.
\]
Il ramo e' inizialmente stabile dopo la transizione,
ma puo' perdere e poi recuperare la stabilita'.
\medskip
Per questo la violazione del criterio di Seidov \emph{non e'}
necessaria perche' esistano stelle gemelle;
e' necessaria solo per l'instabilita' \emph{immediata}.
\end{frame}

\begin{frame}{Perche' conta la pressione di transizione?}
\small
\begin{itemize}
\item Se $P_{\rm trans}$ e' molto alta, la nuova fase appare
vicino al limite massimo del ramo adronico.
\item Potrebbe non esserci spazio sufficiente per una successiva
stabilizzazione della nuova fase.
\item Se la transizione avviene a pressioni minori, il core denso
puo' crescere maggiormente prima della massa massima.
\end{itemize}
\medskip
La retta di Seidov e' locale; gli altri confini della
Figura 9.7 dipendono dalla EOS adronica e dal parametro $s$.
\end{frame}

\begin{frame}{Twin stars e conservazione dei barioni}
\small
Due stelle gemelle possono avere
\[
M_{G,1}\simeq M_{G,2},\qquad R_1
eq R_2.
\]
Ma la massa barionica e'
\[
M_B=m_bN_B,
\]
e in generale
\[
N_{B,1}
eq N_{B,2}.
\]
Per una \emph{trasformazione} fisica a barioni conservati occorre
confrontare configurazioni con lo stesso $N_B$, non semplicemente
con la stessa massa gravitazionale.
\end{frame}

\begin{frame}{Riepilogo: dal salto di fase alle stelle gemelle}
\small
\begin{columns}[T,onlytextwidth]
\begin{column}{0.49\textwidth}
\textbf{Instabilita' iniziale}
\[
\frac{\Delta\varepsilon}{\varepsilon_{\rm trans}}>
\frac12+\frac32\frac{P_{\rm trans}}
{\varepsilon_{\rm trans}}.
\]
\textbf{Fase ad alta densita'}
\[
P=s(\varepsilon-\varepsilon_*).
\]
\end{column}
\begin{column}{0.49\textwidth}
\begin{itemize}
\item $\Delta\varepsilon$: salto energetico.
\item $P_{\rm trans}$: posizione della transizione.
\item $s$: rigidita' della nuova fase.
\item A: nessun ramo ibrido.
\item B: connesso + disconnesso.
\item C: solo connesso.
\item D: solo disconnesso.
\end{itemize}
\end{column}
\end{columns}
\end{frame}

\section{9.4 Mass--Radius Relation with a Phase Transition}

% 01
\begin{frame}{Domanda del paragrafo 9.4}
\small
\begin{columns}[T,onlytextwidth]
\begin{column}{0.49\textwidth}
\textbf{Paragrafo 9.1}
\medskip
EOS nucleonica a basse densita' e pQCD ad alte densita', collegate attraverso interpolazioni lisce.
\[
P=P(\varepsilon).
\]
Vincoli: causalita' e possibilita' di sostenere stelle con $M\gtrsim2M_\odot$.
\end{column}
\begin{column}{0.49\textwidth}
\textbf{Paragrafi 9.3--9.4}
\medskip
Ammettiamo una transizione del primo ordine con salto in densita' di energia.
\[
\varepsilon_H\longrightarrow\varepsilon_Q
\quad\text{a}\quad P=P_{\rm trans}.
\]
Quanto si amplia la regione ammissibile del diagramma massa--raggio?
\end{column}
\end{columns}
\medskip
Il risultato e' l'inviluppo delle curve $M(R)$ ottenute variando l'EOS, non una singola relazione massa--raggio.
\end{frame}

% 02
\begin{frame}{Collegamento con il criterio di Seidov}
\small
\begin{columns}[T,onlytextwidth]
\begin{column}{0.49\textwidth}
\textbf{Transizione di Maxwell}
\[
\begin{gathered}
P(\varepsilon_1)=P(\varepsilon_2)=P_{\rm trans},\\
\varepsilon_1=\varepsilon_{\rm trans},\\
\varepsilon_2=\varepsilon_{\rm trans}+\Delta\varepsilon.
\end{gathered}
\]
Il salto della densita' di energia puo' destabilizzare la configurazione quando nasce un core infinitesimo.
\end{column}
\begin{column}{0.49\textwidth}
\textbf{Instabilita' immediata}
\[
\frac{\Delta\varepsilon}{\varepsilon_{\rm trans}}>
\frac12+\frac32\frac{P_{\rm trans}}{\varepsilon_{\rm trans}}.
\]
E' il criterio di Seidov.
\medskip
\textit{Non garantisce} un nuovo ramo stabile a pressioni centrali maggiori: serve anche una nuova fase sufficientemente rigida.
\end{column}
\end{columns}
\end{frame}

% 03
\begin{frame}{Sofficita' e rigidita' dell'equazione di stato}
\small
La risposta locale della pressione a un incremento di energia e' misurata dalla velocita' del suono:
\[
c_s^2=\frac{dP}{d\varepsilon}\qquad(c=1).
\]
\begin{columns}[T,onlytextwidth]
\begin{column}{0.49\textwidth}
\textbf{Forte ammorbidimento}
\[
c_s^2\simeq0.
\]
Un incremento di $\varepsilon$ produce una variazione molto piccola della pressione.
\medskip
Nel plateau ideale di Maxwell: $\Delta P=0$ per il salto $\Delta\varepsilon$.
\end{column}
\begin{column}{0.49\textwidth}
\textbf{Forte irrigidimento}
\[
c_s^2\longrightarrow1.
\]
La pressione cresce con la massima pendenza consentita dalla condizione causale usata nel testo.
\medskip
Una fase rigida puo' sostenere nuovamente configurazioni compatte.
\end{column}
\end{columns}
\end{frame}

% 04
\begin{frame}{Un salto netto non e' l'unico meccanismo}
\small
\begin{columns}[T,onlytextwidth]
\begin{column}{0.49\textwidth}
\textbf{Maxwell}
\[
P=P_{\rm trans},\qquad
\Delta\varepsilon>0.
\]
La discontinuita' si realizza idealmente in un'interfaccia radiale netta nella stella.
\medskip
Non esiste uno strato TOV macroscopico a pressione costante.
\end{column}
\begin{column}{0.49\textwidth}
\textbf{Gibbs e altre EOS}
\medskip
Una fase mista puo' produrre un ammorbidimento continuo, seguito da un rapido irrigidimento.
\medskip
Anche in questo caso possono comparire rami stabili disconnessi.
\end{column}
\end{columns}
\medskip
Per la terza famiglia conta la struttura globale di $P(\varepsilon)$: un ammorbidimento marcato seguito da un irrigidimento sufficientemente forte.
\end{frame}

\subsection{Vincoli ed EOS estremali}
% 05
\begin{frame}{Vincoli adottati nella costruzione}
\small
\begin{columns}[T,onlytextwidth]
\begin{column}{0.49\textwidth}
\textbf{Materia a basse densita'}
\medskip
Si usa un'EOS nucleonica $P_{\rm low}(\varepsilon)$ vincolata nell'intervallo affidabile.
\medskip
\textbf{Massa massima}
\[
M_{\max}\geq2M_\odot.
\]
La EOS deve sostenere almeno le stelle massicce del riferimento osservativo considerato nel libro.
\end{column}
\begin{column}{0.49\textwidth}
\textbf{Causalita' e stabilita'}
\[
0\leq\frac{dP}{d\varepsilon}\leq1.
\]
La risposta della pressione e' limitata superiormente.
\medskip
\textbf{Liberta' aggiuntiva}
\medskip
Oltre il tratto nucleonico noto si ammettono cambiamenti improvvisi di rigidita' e salti di $\varepsilon$.
\end{column}
\end{columns}
\end{frame}

% 06
\begin{frame}{Due problemi estremali distinti}
\small
\begin{columns}[T,onlytextwidth]
\begin{column}{0.49\textwidth}
\textbf{Massimizzare la massa}
\medskip
Quale prosecuzione ad alta densita' consente il piu' grande valore di $M_{\max}$?
\[
\frac{dP}{d\varepsilon}=1.
\]
Serve la prosecuzione piu' rigida consentita.
\end{column}
\begin{column}{0.49\textwidth}
\textbf{Minimizzare il raggio}
\medskip
Quale EOS produce il raggio piu' piccolo a massa fissata, mantenendo $M_{\max}\geq2M_\odot$?
\[
\underbrace{c_s^2=0}_{\text{salto}}
\ \longrightarrow\
\underbrace{c_s^2=1}_{\text{fase rigida}}.
\]
\end{column}
\end{columns}
\medskip
Le due costruzioni delimitano regioni diverse della Figura 9.8.
\end{frame}

% 07
\begin{frame}{EOS causalmente piu' rigida: integrazione}
\small
Per densita' superiori a una soglia fiduciale $\varepsilon_f$, imponiamo
\[
\frac{dP}{d\varepsilon}=1.
\]
Integrando rispetto a $\varepsilon$:
\[
P(\varepsilon)=\varepsilon+C.
\]
L'EOS deve raccordarsi al valore noto $P_f=P_{\rm low}(\varepsilon_f)$:
\[
P_f=\varepsilon_f+C
\quad\Longrightarrow\quad C=P_f-\varepsilon_f.
\]
Quindi
\[
P(\varepsilon)=P_f+(\varepsilon-\varepsilon_f),
\qquad\varepsilon\geq\varepsilon_f.
\]
\end{frame}

% 08
\begin{frame}{Raccordo alla materia nucleonica}
\small
La EOS completa della costruzione massimamente rigida e'
\[
P(\varepsilon)=
\begin{cases}
P_{\rm low}(\varepsilon),&\varepsilon\leq\varepsilon_f,\\[1mm]
P_f+\varepsilon-\varepsilon_f,&\varepsilon>\varepsilon_f.
\end{cases}
\]
\begin{columns}[T,onlytextwidth]
\begin{column}{0.49\textwidth}
La pressione e' continua al raccordo:
\[
P(\varepsilon_f^-)=P(\varepsilon_f^+)=P_f.
\]
\end{column}
\begin{column}{0.49\textwidth}
La derivata puo' cambiare bruscamente:
\[
\left.\frac{dP}{d\varepsilon}\right|_+=1.
\]
Un salto di pendenza non implica necessariamente un salto di energia.
\end{column}
\end{columns}
\end{frame}

% 09
\begin{frame}{Limite di Rhoades--Ruffini: interpretazione}
\small
\begin{columns}[T,onlytextwidth]
\begin{column}{0.49\textwidth}
Le TOV sono integrate usando
\[
\frac{dm}{dr}=4\pi r^2\varepsilon,
\]
\[
\frac{dP}{dr}
=-\frac{(\varepsilon+P)(m+4\pi r^3P)}{r(r-2m)}
\]
(in unita' $G=c=1$).
\end{column}
\begin{column}{0.49\textwidth}
L'estremo $dP/d\varepsilon=1$ massimizza la risposta pressoria consentita dalla causalita'.
\medskip
Si ottiene un limite superiore alla massa massima \textit{all'interno delle ipotesi adottate sull'EOS nota e sulla prosecuzione}.
\medskip
Non equivale a una previsione che la materia reale abbia $c_s=c$.
\end{column}
\end{columns}
\end{frame}

% 10
\begin{frame}{Risultato della prosecuzione piu' rigida}
\small
Nel modello e con l'EOS nucleonica esterna scelti nel testo:
\[
M_{\max}\simeq3.9M_\odot,\qquad
R(M_{\max})\simeq16.7\,\mathrm{km}.
\]
\begin{columns}[T,onlytextwidth]
\begin{column}{0.49\textwidth}
\textbf{Confronto}
\[
M_{\max}^{\rm RR}\sim4.2M_\odot
\]
per la scelta idealizzata richiamata dal libro, con pressione fiduciale nulla.
\end{column}
\begin{column}{0.49\textwidth}
La differenza e' legata al trattamento dello strato esterno di materia neutronica e al punto fiduciale.
\medskip
Il limite dipende da tali ipotesi: non e' una massa massima universale della QCD.
\end{column}
\end{columns}
\end{frame}

\subsection{Limite inferiore ai raggi}
% 11
\begin{frame}{Perche' un salto energetico puo' ridurre il raggio}
\small
Per una stella di massa gravitazionale $M$ e raggio $R$, la densita' media di massa e'
\[
\bar\rho=\frac{3M}{4\pi R^3}.
\]
A massa fissata,
\[
R\downarrow\quad\Longrightarrow\quad\bar\rho\uparrow.
\]
\begin{columns}[T,onlytextwidth]
\begin{column}{0.49\textwidth}
Un salto di $\varepsilon$ a pressione costante consente un aumento di energia per unita' di volume senza aumento di $P$.
\end{column}
\begin{column}{0.49\textwidth}
Ma un salto troppo grande puo' destabilizzare la stella: occorre un'alta rigidita' nella fase successiva.
\end{column}
\end{columns}
\end{frame}

% 12
\begin{frame}{Primo tratto: EOS nucleonica affidabile}
\small
Fino alla densita' di transizione manteniamo la descrizione nucleonica:
\[
P=P_{\rm low}(\varepsilon),
\qquad\varepsilon<\varepsilon_{\rm trans}.
\]
Definiamo
\[
P_{\rm trans}=P_{\rm low}(\varepsilon_{\rm trans}).
\]
\medskip
Il testo sceglie la transizione alla pressione piu' bassa consentita dal raccordo con l'EOS della materia neutronica utilizzata per costruire il limite estremo.
\medskip
$P_{\rm trans}$ e $\varepsilon_{\rm trans}$ sono i valori \textit{appena prima} dell'apparizione della nuova fase.
\end{frame}

% 13
\begin{frame}{Secondo tratto: transizione di Maxwell}
\small
Alla pressione $P_{\rm trans}$ la densita' passa da
\[
\varepsilon_1=\varepsilon_{\rm trans}
\]
a
\[
\varepsilon_2=\varepsilon_{\rm trans}+\Delta\varepsilon.
\]
La pressione rimane invariata durante la coesistenza ideale:
\[
P(\varepsilon_1)=P(\varepsilon_2)=P_{\rm trans}.
\]
Il plateau nel diagramma $P(\varepsilon)$ codifica il salto:
\[
\Delta P=0,\qquad\Delta\varepsilon>0.
\]
Nelle TOV la densita' cambia a un raggio critico, senza uno strato di pressione perfettamente costante.
\end{frame}

% 14
\begin{frame}{Terzo tratto: massima rigidita' causale}
\small
Dopo la transizione assumiamo
\[
\frac{dP}{d\varepsilon}=1,
\qquad \varepsilon>\varepsilon_2.
\]
Integrare fornisce $P=\varepsilon+C$. Il raccordo al punto $(\varepsilon_2,P_{\rm trans})$ implica
\[
C=P_{\rm trans}-\varepsilon_2.
\]
Quindi
\[
\begin{aligned}
P(\varepsilon)
&=P_{\rm trans}+(\varepsilon-\varepsilon_2)\\
&=P_{\rm trans}+\varepsilon
  -(\varepsilon_{\rm trans}+\Delta\varepsilon).
\end{aligned}
\]
Questa nuova fase puo' ristabilizzare configurazioni prima instabili.
\end{frame}

% 15
\begin{frame}{EOS estrema completa per raggi piccoli}
\footnotesize
\[
P(\varepsilon)=
\begin{cases}
P_{\rm low}(\varepsilon),
  &\varepsilon<\varepsilon_{\rm trans},\\[1.5mm]
P_{\rm trans},
  &\varepsilon_{\rm trans}\leq\varepsilon\leq
  \varepsilon_{\rm trans}+\Delta\varepsilon,\\[1.5mm]
P_{\rm trans}+\varepsilon-
  (\varepsilon_{\rm trans}+\Delta\varepsilon),
  &\varepsilon>\varepsilon_{\rm trans}+\Delta\varepsilon.
\end{cases}
\]
\medskip
\begin{columns}[T,onlytextwidth]
\begin{column}{0.49\textwidth}
\textbf{Tratto iniziale}\par\smallskip
EOS di materia neutronica nota.
\medskip
\textbf{Tratto intermedio}\par\smallskip
Salto di Maxwell, $\Delta\varepsilon>0$ a $P_{\rm trans}$.
\end{column}
\begin{column}{0.49\textwidth}
\textbf{Tratto finale}\par\smallskip
EOS causalmente piu' rigida, $c_s^2=1$.
\medskip
E' una costruzione estremale, non una descrizione microscopica univoca.
\end{column}
\end{columns}
\end{frame}

% 16
\begin{frame}{Forma inversa dell'EOS per le TOV}
\small
Nell'integrazione numerica e' spesso piu' comodo usare $\varepsilon=\varepsilon(P)$.
\medskip
Per la fase ad alta densita':
\[
\begin{aligned}
P-P_{\rm trans}
&=\varepsilon-(\varepsilon_{\rm trans}+\Delta\varepsilon),\\
\varepsilon_Q(P)
&=\varepsilon_{\rm trans}+\Delta\varepsilon
  +(P-P_{\rm trans}).
\end{aligned}
\]
Per la fase a bassa densita':
\[
\varepsilon_H(P)=\varepsilon_{\rm low}(P).
\]
Al passaggio per $P=P_{\rm trans}$, la densita' cambia da $\varepsilon_Q$ a $\varepsilon_H$; $m(r)$ e $P(r)$ rimangono continui.
\end{frame}

% 17
\begin{frame}{Schema dell'EOS estrema nel piano $P$--$\varepsilon$}
\centering
\begin{tikzpicture}[x=0.96cm,y=0.87cm]
\draw[->,thick] (0,0)--(10.9,0) node[below] {$\varepsilon$};
\draw[->,thick] (0,0)--(0,6.0) node[left] {$P$};
\draw[thick] (0.3,0.45) .. controls (1.8,0.55) and (2.8,1.2) .. (4.1,2.65);
\draw[thick] (4.1,2.65)--(6.55,2.65);
\draw[thick] (6.55,2.65)--(8.85,5.55);
\draw[densely dashed,thin] (4.1,0)--(4.1,2.65);
\draw[densely dashed,thin] (6.55,0)--(6.55,2.65);
\node[below] at (4.1,0) {$\varepsilon_1$};
\node[below] at (6.55,0) {$\varepsilon_2$};
\node[left] at (0,2.65) {$P_{\rm trans}$};
\draw[dotted] (0,2.65)--(4.1,2.65);
\node[align=center,font=\scriptsize] at (2.25,2.85) {EOS nucleonica};
\node[font=\scriptsize,above] at (5.4,2.65) {$\Delta\varepsilon$};
\node[align=center,font=\scriptsize] at (8.75,4.95) {$c_s^2=1$};
\node[font=\scriptsize] at (5.2,-0.85) {Schema qualitativo, non in scala};
\end{tikzpicture}
\end{frame}

% 18
\begin{frame}{Due effetti opposti nella stessa EOS}
\small
\begin{columns}[T,onlytextwidth]
\begin{column}{0.49\textwidth}
\textbf{Salto $\Delta\varepsilon$}
\[
P=\text{costante},\quad\Delta\varepsilon>0.
\]
Aumenta la densita' di energia senza un aumento della pressione.
\medskip
Favorisce strutture molto compatte, ma puo' innescare instabilita' (Seidov).
\end{column}
\begin{column}{0.49\textwidth}
\textbf{Fase rigida dopo il salto}
\[
\frac{dP}{d\varepsilon}=1.
\]
Fornisce il massimo incremento pressorio consentito dalla causalita'.
\medskip
Puo' sostenere un nuovo ramo stabile e aumentare la massa massima.
\end{column}
\end{columns}
\medskip
Il raggio piu' piccolo ammissibile deriva dall'\emph{ottimizzazione congiunta} dei due effetti, non da un salto arbitrariamente grande.
\end{frame}

% 19
\begin{frame}{Come si sceglie il salto energetico}
\small
Il testo impone una configurazione di massa massima non inferiore al riferimento osservativo:
\[
M_{\max}\geq2M_\odot.
\]
Per costruire il bordo di raggi minimi si regola $\Delta\varepsilon$ fino al caso estremo
\[
M_{\max}=2M_\odot.
\]
\begin{columns}[T,onlytextwidth]
\begin{column}{0.49\textwidth}
Se il salto e' troppo debole, si ottengono in generale raggi maggiori del minimo cercato.
\end{column}
\begin{column}{0.49\textwidth}
Se e' troppo forte, la sequenza rischia di non sostenere piu' $2M_\odot$.
\end{column}
\end{columns}
\medskip
Scelta addizionale: la transizione avviene al punto piu' basso ammesso dalla EOS nucleonica di riferimento.
\end{frame}

\subsection{Figura 9.8 e limiti massa--raggio}
% 20
\begin{frame}{Figura 9.8: massa--raggio con transizione}
    \begin{figure}[!ht]\includegraphics[trim={0cm 0cm 0 0},clip,keepaspectratio,width=0.95\textwidth]{Schaffner/MR-hybrid-phasetransition}\label{fig:Schaffner_MR-hybrid-phasetransition}\end{figure}
\end{frame}

% 21
\begin{frame}{Come leggere le due regioni della Figura 9.8}
\small
\begin{columns}[T,onlytextwidth]
\begin{column}{0.49\textwidth}
\textbf{Regione chiara}
\medskip
Famiglia di EOS interpolate in modo liscio fra il vincolo nucleonico e la pQCD, come nella Figura 9.2.
\medskip
Corrisponde alla classe di parametrizzazioni senza salti netti considerata nel \S9.1.
\end{column}
\begin{column}{0.49\textwidth}
\textbf{Regione scura}
\medskip
Insieme piu' ampio ottenuto ammettendo forti variazioni della rigidita' e transizioni con salto energetico.
\medskip
Include raggi notevolmente inferiori a quelli della regione chiara.
\end{column}
\end{columns}
\medskip
Le aree non sono singole sequenze TOV: ogni EOS genera una propria curva $M(R)$; l'unione di molte curve determina una regione.
\end{frame}

% 22
\begin{frame}{Una differenza di ipotesi, non un singolo modello}
\small
\begin{columns}[T,onlytextwidth]
\begin{column}{0.49\textwidth}
\textbf{Figura 9.2}
\[
P(\varepsilon)\ \text{raccordata senza salti}.
\]
Sono utilizzate politropiche a tratti e condizioni di matching con una regione di pQCD.
\end{column}
\begin{column}{0.49\textwidth}
\textbf{Figura 9.8}
\[
\Delta\varepsilon>0
\quad\text{a}\quad P_{\rm trans}.
\]
Vengono ammesse anche prosecuzioni causali estreme dopo la transizione.
\end{column}
\end{columns}
\medskip
Le due fasce non differiscono \emph{solo} per la possibilita' di un salto: cambiano anche le famiglie di parametrizzazioni utilizzate. La regione scura e' l'insieme piu' ampio nelle ipotesi degli autori.
\end{frame}

% 23
\begin{frame}{Rette di compattezza nel diagramma $M$--$R$}
\small
Definiamo
\[
\mathcal C=\frac{GM}{Rc^2}.
\]
A compattezza fissata:
\[
M=\mathcal C\frac{c^2}{G}R.
\]
\begin{columns}[T,onlytextwidth]
\begin{column}{0.49\textwidth}
Sono rette passanti per l'origine nel piano $M$--$R$.
\medskip
La Figura 9.8 mostra anche $\mathcal C=0.1$, $0.2$ e $0.3$.
\end{column}
\begin{column}{0.49\textwidth}
Raggio piu' piccolo a massa fissata significa compattezza maggiore:
\[
R\downarrow\quad\Longrightarrow\quad\mathcal C\uparrow.
\]
\end{column}
\end{columns}
\end{frame}

% 24
\begin{frame}{Tre limiti relativistici riportati nella figura}
\small
\begin{align*}
\text{Schwarzschild:}&\qquad
R>\frac{2GM}{c^2},&\mathcal C&<\frac12,\\[1mm]
\text{Buchdahl:}&\qquad
R\geq\frac94\frac{GM}{c^2},&\mathcal C&\leq\frac49,\\[1mm]
\text{Causalita':}&\qquad
R\gtrsim\frac{GM}{0.354\,c^2},&\mathcal C&\lesssim0.354.
\end{align*}
\medskip
\begin{columns}[T,onlytextwidth]
\begin{column}{0.49\textwidth}
Il limite di Buchdahl richiede ipotesi specifiche su simmetria, isotropia e profilo di densita'.
\end{column}
\begin{column}{0.49\textwidth}
Il valore $0.354$ e' un limite globale di compattezza per la classe di EOS causali e stelle stabili considerata: non deriva soltanto dalla disuguaglianza locale $c_s\leq c$.
\end{column}
\end{columns}
\end{frame}

% 25
\begin{frame}{Limite causale al raggio per $M=2M_\odot$}
\small
Dalla compattezza massima del libro:
\[
\mathcal C\leq0.354
\quad\Longrightarrow\quad
R\geq\frac{GM}{0.354\,c^2}.
\]
Usiamo la scala gravitazionale solare:
\[
\frac{GM_\odot}{c^2}=1.4766\,\mathrm{km}.
\]
Per $M=2M_\odot$:
\[
\begin{aligned}
R_{\rm caus}
&=\frac{2\times1.4766\,\mathrm{km}}{0.354}\\
&\simeq8.34\,\mathrm{km}.
\end{aligned}
\]
\end{frame}

% 26
\begin{frame}{Raggio minimo con transizione: confronto}
\small
Il bordo sinistro calcolato nella Figura 9.8 raggiunge, per la configurazione limite,
\[
M_{\max}=2M_\odot,\qquad
R_{\min}^{\rm hybrid}\simeq8.6\,\mathrm{km}.
\]
\medskip
Confrontiamo con il limite causale:
\[
R_{\rm caus}\simeq8.3\,\mathrm{km}.
\]
La differenza e'
\[
\Delta R\simeq8.6-8.3\simeq0.3\,\mathrm{km}.
\]
La costruzione estrema con transizione e successiva fase rigida si avvicina molto al limite di compattezza della classe considerata.
\end{frame}

% 27
\begin{frame}{Perche' le curve limite si avvicinano alle alte masse}
\small
\begin{columns}[T,onlytextwidth]
\begin{column}{0.49\textwidth}
Nelle stelle molto massicce il core ad alta densita' ha un peso maggiore nella struttura globale.
\medskip
La parte esterna di materia neutronica influisce relativamente meno sulla configurazione.
\end{column}
\begin{column}{0.49\textwidth}
La curva del raggio minimo puo' quindi avvicinarsi a quella dell'EOS causale ideale.
\medskip
Il testo confronta questo comportamento con stelle strange con e senza crosta (Figura 8.4).
\end{column}
\end{columns}
\medskip
Le proprieta' degli strati esterni rimangono importanti soprattutto per la geometria della superficie e le configurazioni meno massicce.
\end{frame}

% 28
\begin{frame}{L'altro estremo della regione: masse molto grandi}
\small
La prosecuzione massimamente rigida senza il forte plateau utilizzato per minimizzare $R$ da' il bordo di grande massa della regione.
\[
M_{\max}\simeq3.9M_\odot,
\qquad R\simeq16.7\,\mathrm{km}.
\]
Confronto con il limite idealizzato citato:
\[
M_{\max}^{\rm RR}\simeq4.2M_\odot.
\]
Il risultato dipende dall'EOS nota a basse densita' e dalla scelta fiduciale.
\medskip
Non e' una previsione della massa massima realmente raggiunta dalla materia stellare: e' un vincolo estremale sotto ipotesi specifiche.
\end{frame}

\subsection{Interpretazione fisica dei vincoli}
% 29
\begin{frame}{Quale conclusione osservativa permette la Figura 9.8?}
\small
Supponiamo che una misura futura fornisca
\[
(M_{\rm obs},R_{\rm obs})
\]
al di fuori della banda chiara ma dentro la regione scura.
\medskip
\begin{columns}[T,onlytextwidth]
\begin{column}{0.49\textwidth}
\textbf{Conclusione diretta}
\medskip
La classe di EOS lisce utilizzata per la banda chiara non descrive quella configurazione, nei limiti delle incertezze.
\end{column}
\begin{column}{0.49\textwidth}
\textbf{Interpretazione possibile}
\medskip
La vera EOS potrebbe presentare un forte cambiamento di rigidita' oppure una transizione del primo ordine.
\end{column}
\end{columns}
\medskip
Questo non identifica automaticamente la nuova fase con materia di quark: la struttura gravitazionale dipende da $P(\varepsilon)$.
\end{frame}

% 30
\begin{frame}{Terza famiglia e regione ammissibile: due domande diverse}
\small
\begin{columns}[T,onlytextwidth]
\begin{column}{0.49\textwidth}
\textbf{Paragrafo 9.3}
\medskip
Per una EOS fissata, quanti rami stabili delle TOV esistono?
\[
\text{stabile}\to\text{instabile}
\to\text{stabile}.
\]
Possibili twin stars con masse simili e raggi differenti.
\end{column}
\begin{column}{0.49\textwidth}
\textbf{Paragrafo 9.4}
\medskip
Variando le EOS ammissibili, quali sono gli estremi raggiungibili nel diagramma $M$--$R$?
\[
R_{\min}(M),\qquad M_{\max}^{\rm upper}.
\]
L'inviluppo non dimostra da solo l'esistenza di rami disconnessi.
\end{column}
\end{columns}
\end{frame}

% 31
\begin{frame}{Riepilogo: equazioni e numeri del paragrafo}
\footnotesize
\begin{columns}[T,onlytextwidth]
\begin{column}{0.49\textwidth}
\textbf{Causalita'}
\[
0\leq\frac{dP}{d\varepsilon}\leq1.
\]
\textbf{Transizione di Maxwell}
\[
P=P_{\rm trans},\quad\Delta\varepsilon>0.
\]
\textbf{Fase estremamente rigida}
\[
P=P_{\rm trans}+\varepsilon-
(\varepsilon_{\rm trans}+\Delta\varepsilon).
\]
\end{column}
\begin{column}{0.49\textwidth}
\textbf{Vincolo osservativo adottato}
\[
M_{\max}\geq2M_\odot.
\]
\textbf{Bordo di raggio minimo}
\[
R(2M_\odot)\simeq8.6\,\mathrm{km}.
\]
\textbf{Bordo di massa elevata}
\[
M_{\max}\simeq3.9M_\odot,
\quad R\simeq16.7\,\mathrm{km}.
\]
\end{column}
\end{columns}
\end{frame}

% 32
\begin{frame}{Conclusione del paragrafo 9.4}
\small
\begin{columns}[T,onlytextwidth]
\begin{column}{0.49\textwidth}
Una EOS con forte ammorbidimento e successivo irrigidimento puo' produrre stelle piu' compatte di quanto prevedano le interpolazioni lisce considerate.
\medskip
L'ammissione di transizioni di fase amplia la regione ammissibile del piano $M$--$R$.
\end{column}
\begin{column}{0.49\textwidth}
Osservazioni precise di masse e raggi possono vincolare forti cambiamenti nell'EOS.
\medskip
La composizione microscopica non e' determinata univocamente: EOS diverse possono produrre curve $M(R)$ simili.
\end{column}
\end{columns}
\medskip
Il confronto con le osservazioni deve distinguere il vincolo sull'andamento di $P(\varepsilon)$ dall'interpretazione della possibile transizione adroni--quark.
\end{frame}

